% Apostila de Programacao Inteira (Modelagem) - GCC118 Programacao Matematica (UFLA)
% Guia de estudo para o notebook Aula_Programacao_Inteira_Modelagem.ipynb
% Compilar a partir desta pasta:  tectonic apostila_programacao_inteira.tex   (ou xelatex/lualatex/pdflatex)
\documentclass[11pt,a4paper,oneside]{report}

\usepackage{iftex}
\ifPDFTeX
  \usepackage[utf8]{inputenc}
  \usepackage[T1]{fontenc}
  \usepackage{lmodern}
\else
  \usepackage{fontspec}
\fi
\usepackage[brazil]{babel}
\usepackage[margin=2.4cm]{geometry}
\usepackage{amsmath,amssymb,amsthm}
\usepackage{icomma}
\usepackage{booktabs,tabularx,array,multirow}
\usepackage{enumitem}
\usepackage{xcolor}
\usepackage{ifthen}
\usepackage{tikz}
\usetikzlibrary{arrows.meta,positioning,patterns,shapes.geometric,calc,decorations.pathreplacing}
\usepackage[most]{tcolorbox}
\usepackage{listings}
\usepackage{caption}
\usepackage{float}
\usepackage[hidelinks]{hyperref}
\usepackage{xurl}

\captionsetup{font=small,labelfont=bf}
\setlist{leftmargin=*,itemsep=2pt,topsep=4pt}
\setlength{\parskip}{4pt}
\setlength{\emergencystretch}{3em}
\raggedbottom
\renewcommand{\arraystretch}{1.15}

% ---------------------------------------------------------------- cores
\definecolor{azul}{HTML}{1F4E79}
\definecolor{verde}{HTML}{2E7D32}
\definecolor{laranja}{HTML}{E65100}
\definecolor{roxo}{HTML}{5E35B1}
\definecolor{cinza}{HTML}{555555}
\definecolor{fundocodigo}{HTML}{F6F8FA}

% ---------------------------------------------------------------- caixas
\tcbset{base/.style={enhanced,breakable,boxrule=0.6pt,arc=2pt,left=6pt,right=6pt,top=4pt,bottom=4pt,
  fonttitle=\bfseries\small,before skip=8pt,after skip=8pt}}
\newtcolorbox{definicao}[1][]{base,colback=azul!4,colframe=azul,title={Definição\ifx&#1&\else: #1\fi}}
\newtcolorbox{exemplo}[1][]{base,colback=verde!4,colframe=verde,title={Exemplo prático\ifx&#1&\else: #1\fi}}
\newtcolorbox{atencao}[1][]{base,colback=laranja!5,colframe=laranja,title={Cuidado\ifx&#1&\else: #1\fi}}
\newtcolorbox{receita}[1][]{base,colback=roxo!4,colframe=roxo,title={Receita\ifx&#1&\else: #1\fi}}
\newtcolorbox{resumo}[1][]{base,colback=cinza!5,colframe=cinza,title={\ifx&#1&Em resumo\else #1\fi}}
\newtcolorbox{discussao}[1][]{base,colback=yellow!6,colframe=yellow!50!black,title={Perguntas da discussão\ifx&#1&\else: #1\fi}}
\newtcolorbox{outros}[1][]{base,colback=teal!4,colframe=teal!60!black,title={Outros autores\ifx&#1&\else: #1\fi}}

% ---------------------------------------------------------------- codigo
\lstdefinestyle{py}{
  language=Python,
  basicstyle=\ttfamily\small,
  keywordstyle=\color{azul}\bfseries,
  commentstyle=\color{verde!70!black}\itshape,
  stringstyle=\color{laranja!80!black},
  columns=fullflexible, keepspaces=true,
  breaklines=true, showstringspaces=false,
  aboveskip=0pt, belowskip=0pt,
  morekeywords={True,False,None,as,with},
  extendedchars=true,
  literate={á}{{\'a}}1 {à}{{\`a}}1 {ã}{{\~a}}1 {â}{{\^a}}1 {é}{{\'e}}1 {ê}{{\^e}}1
           {í}{{\'i}}1 {ó}{{\'o}}1 {õ}{{\~o}}1 {ô}{{\^o}}1 {ú}{{\'u}}1 {ç}{{\c c}}1
           {Á}{{\'A}}1 {É}{{\'E}}1 {Í}{{\'I}}1 {Ó}{{\'O}}1 {Ú}{{\'U}}1 {Ç}{{\c C}}1
}
\lstset{style=py}
\newtcblisting{codigo}{enhanced,listing only,listing options={style=py},
  colback=fundocodigo,colframe=black!15,boxrule=0.5pt,arc=1pt,left=4pt,right=4pt,top=2pt,bottom=2pt,
  before skip=6pt,after skip=6pt}
\newcommand{\py}[1]{\texttt{\detokenize{#1}}}

% ---------------------------------------------------------------- matematica
\newtheorem{teorema}{Teorema}[chapter]
\newcommand{\R}{\mathbb{R}}
\newcommand{\Z}{\mathbb{Z}}
\newcommand{\B}{\{0,1\}}
\newcommand{\xor}{\mathbin{\triangledown}}
\newcommand{\sa}{\text{s.a.}\quad}

\title{%
  {\Huge\bfseries Programação Inteira}\\[6pt]
  {\Large Modelagem com variáveis inteiras e binárias}\\[14pt]
  {\large Apostila de estudo para o notebook\\ \texttt{Aula\_Programacao\_Inteira\_Modelagem.ipynb}}}
\author{GCC118 -- Programação Matemática\\
  \small Universidade Federal de Lavras (UFLA), semestre 2026/2\\[4pt]
  \small Material de apoio elaborado a partir das aulas da Profa.\ Andreza C.\ Beezão Moreira\\
  \small e do Prof.\ Mayron César O.\ Moreira e dos slides \emph{Programação inteira -- Modelagem}}
\date{Outubro de 2026}

\begin{document}
\pagenumbering{roman}
\hypersetup{pageanchor=false}
\maketitle
\hypersetup{pageanchor=true}
\tableofcontents

\clearpage
\pagenumbering{arabic}

\chapter*{Como usar esta apostila}
\addcontentsline{toc}{chapter}{Como usar esta apostila}

Esta apostila foi escrita para quem está tendo o \textbf{primeiro contato com programação inteira}. Supomos
apenas que você já viu o básico de programação linear (variáveis, função objetivo, restrições, região viável) e que
sabe programar em Python. Tudo o mais é construído aqui, com calma e com exemplos.

O objetivo é que, ao final, você consiga resolver sozinho(a) todos os exercícios do notebook
\texttt{Aula\_Programacao\_Inteira\_Modelagem.ipynb} e, mais importante, \emph{entenda por que} cada modelo é
escrito do jeito que é.

\section*{Organização}

\begin{description}
  \item[Parte I -- Fundamentos.] O que é programação inteira, por que ela é necessária, por que ``resolver como PL e
  arredondar'' não funciona, o que é relaxação linear, como se formula um modelo, o que é o famoso \emph{big-M}, um
  pouco de lógica com variáveis binárias e um tutorial de PuLP (a biblioteca Python usada no notebook).
  \item[Parte II -- O notebook passo a passo.] Um capítulo por bloco do notebook (blocos 0 a 7). Em cada um:
  o problema, as respostas às perguntas da discussão, a construção do modelo matemático, a tradução para código, a
  interpretação do resultado e os erros mais comuns.
  \item[Apêndice.] Respostas dos exercícios de fixação e tabela com os valores ótimos esperados.
\end{description}

\section*{Como o notebook está organizado}

Cada bloco do notebook segue sempre a mesma dinâmica:
\begin{center}
\textbf{Discussão} $\;\rightarrow\;$ \textbf{Modelo} (no quadro) $\;\rightarrow\;$ \textbf{Exercício} (você completa os
\texttt{\# TODO}) $\;\rightarrow\;$ \textbf{Verificação}
\end{center}
As células de \textbf{Verificação} comparam o seu resultado com o gabarito \emph{sem revelar o valor}: aparece
``correto'' ou ``não bate com o gabarito'' (às vezes com uma dica). As células de \textbf{Solução de referência} ficam
ocultas; rodá-las sobrescreve a sua versão. Sugestão: só abra a solução depois de tentar, e use esta apostila
como ``professor particular'' quando travar.

\section*{Caixas usadas no texto}

\begin{definicao}[conceito formal]
Definições que você precisa saber de cor.
\end{definicao}
\begin{exemplo}[exemplo resolvido]
Pequenos problemas resolvidos do início ao fim, com todos os passos.
\end{exemplo}
\begin{receita}[técnica de modelagem]
Uma ``fórmula'' pronta para transformar uma regra em restrições lineares.
\end{receita}
\begin{atencao}[erros comuns]
Armadilhas em que quase todo mundo cai na primeira vez.
\end{atencao}
\begin{outros}[outras fontes]
Complementos de livros, artigos, cursos abertos e documentação de solvers, sempre com a referência
(lista completa no fim da apostila).
\end{outros}

\section*{Fontes}

A apostila segue o notebook da disciplina~\cite{notebook} e os slides \emph{Programação inteira --
Modelagem}~\cite{slides}, que por sua vez se baseiam em Wolsey~\cite{wolsey1998}, Arenales et
al.~\cite{arenales2007} e Sierksma~\cite{sierksma1996}. Para ampliar a visão, incluímos material de outros autores,
como o capítulo de programação inteira do livro do MIT de Bradley, Hax e Magnanti~\cite{amp}, artigos históricos
(Land e Doig, Gomory, Dantzig, Smith), surveys e a documentação de solvers comerciais. Todas as fontes estão nas
Referências.

\section*{Roteiro de estudo sugerido}

\begin{enumerate}
  \item Leia os Capítulos~1 a 4 antes da aula (ou antes de abrir o notebook).
  \item Faça os exercícios de fixação do fim de cada capítulo e confira no Apêndice.
  \item Leia o Capítulo~5 (PuLP) com o notebook aberto e rode as duas primeiras células.
  \item Para cada bloco do notebook: tente sozinho(a); se travar, leia a seção correspondente da Parte~II até o ponto em
  que destravar, e volte ao código.
  \item No fim, revise o quadro-resumo do Capítulo~14.
\end{enumerate}


\part{Fundamentos de programação inteira}
\chapter{Da programação linear à programação inteira}
\label{cap:pl-pi}

\begin{flushright}
\begin{minipage}{0.75\linewidth}\itshape\small
``A arte de formular programas lineares e lineares inteiros é, bem, uma arte: é difícil de ensinar e ainda mais
difícil de aprender.''
\end{minipage}\\[2pt]
{\small --- G.~G.~Brown e R.~F.~Dell~\cite{browndell2007}, citados em~\cite{ravelo}}
\end{flushright}

A boa notícia é que essa arte tem técnicas que se repetem. O objetivo desta apostila é apresentá-las uma a uma,
com exemplos.

\section{Relembrando: o que é um problema de programação linear}

Um problema de \textbf{programação linear} (PL) tem três ingredientes:
\begin{enumerate}
  \item \textbf{Variáveis de decisão}: os números que queremos escolher (quanto produzir, quanto transportar\ldots).
  \item \textbf{Função objetivo}: uma expressão \emph{linear} nas variáveis que queremos maximizar ou minimizar.
  \item \textbf{Restrições}: desigualdades ou igualdades \emph{lineares} que as variáveis precisam respeitar.
\end{enumerate}
``Linear'' significa que cada variável aparece apenas multiplicada por uma constante e somada às outras: $3x_1 + 2x_2$
é linear; $x_1x_2$, $x_1^2$, $\sqrt{x_1}$ e $\max\{x_1,x_2\}$ não são.

\begin{exemplo}[produção de fósforos (slides)]
Uma empresa produz caixas de fósforos longos (lucro 3 por caixa) e curtos (lucro 2), em unidades de 100\,000 caixas.
Sejam $x_1$ e $x_2$ as quantidades produzidas.
\[
\begin{aligned}
\max\;& 3x_1 + 2x_2\\
\sa & x_1 + x_2 \le 9 &&\text{(capacidade da máquina)}\\
    & 3x_1 + x_2 \le 18 &&\text{(madeira disponível)}\\
    & x_1 \le 7,\quad x_2 \le 6 &&\text{(caixas disponíveis)}\\
    & x_1, x_2 \ge 0.
\end{aligned}
\]
Os pontos que satisfazem todas as restrições formam a \textbf{região viável}, um polígono convexo. O teorema
fundamental da PL garante que, se existe ótimo, existe um ótimo em um \textbf{vértice} desse polígono, e é isso que o
método simplex explora. Os vértices aqui são $(0,0)$, $(6,0)$, $(4{,}5;\,4{,}5)$, $(3,6)$ e $(0,6)$, com valores
$0$, $18$, $22{,}5$, $21$ e $12$. O ótimo é $x_1=x_2=4{,}5$, ou seja, \textbf{450\,000 caixas de cada tipo}.

\emph{O valor fracionário é um problema?} Não: a unidade é ``100\,000 caixas'', então 4,5 significa 450\,000
caixas. Variáveis contínuas são uma \emph{aproximação realista} aqui.
\end{exemplo}

\section{Quando as variáveis precisam ser inteiras}

Agora imagine que a variável é o número de aviões que uma companhia compra, de enfermeiras escaladas para um turno ou
de caminhões enviados a um porto. Não existe ``3,7 aviões''. Land e Doig, no artigo que deu origem ao método
\emph{branch-and-bound}, já diziam que problemas com variáveis inteiras surgem sempre que há
\textbf{indivisibilidades}.

Há dois motivos bem diferentes para usar variáveis inteiras, e vale a pena separá-los desde já:
\begin{enumerate}
  \item \textbf{Natureza do problema}: a quantidade é indivisível (aviões, pessoas, veículos). É o
  ``motivo óbvio''.
  \item \textbf{Artifício de modelagem}: usamos variáveis \emph{binárias} (que só valem 0 ou 1) para representar
  \textbf{decisões} (abrir ou não uma fábrica) e \textbf{lógica} (``se\ldots então'', ``ou'', ``no máximo $k$'',
  custo fixo\ldots) dentro de um modelo que continua linear. É o \textbf{foco do notebook} a partir do Bloco~1.
\end{enumerate}

\begin{outros}[a mesma distinção no livro do MIT]
Bradley, Hax e Magnanti~\cite[Seç.~9.2]{amp} fazem a mesma separação. De um lado, a inteireza ``pela natureza do
problema'', com o exemplo de decidir quantos porta-aviões nucleares uma marinha deve ter: a resposta é da ordem de
um ou dois, e uma fração não teria significado. De outro, variáveis inteiras que surgem de (i) \emph{condições
lógicas} (``se um novo produto for desenvolvido, então uma nova fábrica deve ser construída'') e (ii)
\emph{não linearidades} como custos fixos. Os autores chamam as binárias de variáveis ``sim/não'' (\emph{go--no-go}),
típicas de decisões estratégicas de alto custo, como investimentos.
\end{outros}

\section{As classes de problemas}

\begin{definicao}[tipos de problemas inteiros]
Seja $x$ um vetor de variáveis contínuas e $y$ um vetor de variáveis inteiras.
\begin{itemize}
  \item \textbf{Programação linear inteira mista (PLIM)}: há variáveis contínuas \emph{e} inteiras.
  \[ \max\; cx + dy \quad \sa Ax + Gy \le b,\;\; x \ge 0,\;\; y \ge 0 \text{ e inteiro.}\]
  \item \textbf{Programação linear inteira (PLI)}: todas as variáveis são inteiras.
  \[ \max\; cx \quad \sa Ax \le b,\;\; x \in \Z^n_+.\]
  \item \textbf{Programação binária (PB ou 0--1)}: todas as variáveis são binárias.
  \[ \max\; cx \quad \sa Ax \le b,\;\; x \in \B^n.\]
\end{itemize}
\end{definicao}

Uma família muito importante de problemas é a de \textbf{otimização combinatória}: temos um conjunto finito
$N=\{1,\dots,n\}$ de ``elementos'' (itens, arestas, tarefas\ldots), cada um com um valor $c_j$, e uma família
$\mathcal{F}$ de subconjuntos \emph{viáveis} de $N$. Queremos o melhor subconjunto viável:
\[
\max_{S\subseteq N}\Big\{\sum_{j\in S} c_j \;:\; S\in\mathcal{F}\Big\}.
\]
A mochila (escolher itens que cabem), o caixeiro viajante (escolher arestas que formam um ciclo) e o \emph{bin
packing} são exemplos. Quase sempre eles são escritos como \textbf{programas binários}: $x_j=1$ se $j\in S$.

\begin{atencao}[``finito'' não quer dizer ``fácil'']
Com $n$ itens, há $2^n$ subconjuntos. Para $n=60$, são mais de $10^{18}$ possibilidades: enumerar todas, mesmo a
um bilhão por segundo, levaria mais de 30 anos. Por isso precisamos de bons modelos e de algoritmos como o
\emph{branch-and-bound} (Seção~\ref{sec:bb}).
\end{atencao}

\section{Um primeiro modelo inteiro, do zero}

\begin{exemplo}[aluguel de veículos para uma excursão]
Um centro acadêmico precisa transportar \textbf{50 pessoas}. Pode alugar vans (12 lugares, R\$ 250 cada) e ônibus
(30 lugares, R\$ 600 cada). Quantos de cada para gastar o mínimo?

\textbf{Passo 1 -- variáveis.} $v$ = número de vans e $o$ = número de ônibus. Não se aluga meia van:
$v,o\in\Z_+$.

\textbf{Passo 2 -- restrições.} Os lugares precisam bastar: $12v + 30o \ge 50$.

\textbf{Passo 3 -- objetivo.} $\min\; 250v + 600o$.

\textbf{Passo 4 -- resolver.} Se ignorarmos a inteireza (PL), o ônibus custa $600/30=20$ por lugar e a van
$250/12\approx 20{,}83$; a PL usa só ônibus: $o=50/30\approx1{,}67$, custo $1000$. Como o problema é pequeno,
podemos enumerar as opções inteiras com o mínimo de vans para cada número de ônibus:
\begin{center}
\begin{tabular}{cccc}
\toprule
$o$ & menor $v$ com $12v \ge 50-30o$ & lugares & custo\\
\midrule
0 & 5 & 60 & 1250\\
1 & 2 & 54 & \textbf{1100}\\
2 & 0 & 60 & 1200\\
\bottomrule
\end{tabular}
\end{center}
O ótimo inteiro é \textbf{1 ônibus e 2 vans (R\$ 1100)}. Repare que ``arredondar a PL'' ($o=1{,}67\to2$) dá R\$ 1200:
é viável, mas \emph{não é ótimo}. E o valor da PL (1000) é menor que o ótimo inteiro (1100), porque a PL tem mais
liberdade. Essas duas observações são o tema dos próximos dois capítulos.
\end{exemplo}

\section{Um pouco de história}

Os métodos que o solver do notebook usa têm mais de 60 anos, e vários dos problemas desta apostila são, eles
mesmos, marcos históricos.

\begin{figure}[H]
\centering
\begin{tikzpicture}[x=0.18cm,
  ev/.style={font=\scriptsize,align=center,text width=2.35cm,inner sep=1pt}]
  \draw[thick,-{Stealth}] (0,0) -- (88,0);
  \foreach \ano/\pos/\lado/\txt in {
      6/1939/above/{Kantorovich: planejamento da produção (modelo de corte e empacotamento)},
      14/1947/below/{Dantzig: método simplex},
      23/1956/above/{Smith: regra ótima para uma máquina (WSPT)},
      33/1958/below/{Dantzig (1957): limite da mochila; Gomory (1958): planos de corte},
      44/1960/above/{Land e Doig: \emph{branch-and-bound}},
      54/1965/below/{Dakin: ramificação $x\le\lfloor v\rfloor$ / $x\ge\lceil v\rceil$},
      68/{fim dos anos 1990}/above/{solvers comerciais combinam B\&B e planos de corte (\emph{branch-and-cut})},
      82/2020/below/{2001--2020: solvers de PLIM $\approx$1000$\times$ mais rápidos}}
  { \fill[azul] (\ano,0) circle (2pt);
    \ifthenelse{\equal{\lado}{above}}
      {\draw (\ano,0) -- (\ano,0.55); \node[ev,above] at (\ano,0.55) {\textbf{\pos}\\ \txt};}
      {\draw (\ano,0) -- (\ano,-0.55); \node[ev,below] at (\ano,-0.55) {\textbf{\pos}\\ \txt};} }
\end{tikzpicture}
\caption{Alguns marcos da programação inteira citados nesta apostila. Fontes: Kantorovich~\cite{kantorovich1960};
Smith~\cite{smith1956}; Dantzig~\cite{dantzig1957}; Gomory~\cite{gomory1958}; Land e Doig~\cite{landdoig1960};
Dakin~\cite{dakin1965}; Bixby~\cite{bixby2012}; Koch et al.~\cite{koch2022}. (Posições fora de escala.)}
\label{fig:historia}
\end{figure}

\begin{outros}[Ailsa Land e Alison Doig]
O \emph{branch-and-bound} nasceu de um trabalho contratado pela British Petroleum na London School of Economics, para
melhorar modelos de PL de refinarias em que algumas variáveis precisavam ser inteiras. O artigo de Land e Doig saiu
na \emph{Econometrica} em 1960~\cite{landdoig1960}. Ailsa Land foi a primeira mulher a ser professora titular de
Pesquisa Operacional na Grã-Bretanha~\cite{informsland}. A ideia de resolver relaxações lineares e dividir o problema
em subproblemas ainda é a espinha dorsal dos solvers modernos.
\end{outros}

\begin{resumo}
\begin{itemize}
  \item PL: variáveis contínuas, tudo linear, ótimo em vértice, resolvida eficientemente pelo simplex.
  \item PI: algumas (PLIM) ou todas (PLI) as variáveis são inteiras; PB quando são todas binárias.
  \item Inteireza aparece por \emph{natureza} (indivisibilidade) ou como \emph{ferramenta de modelagem} (binárias
  para decisões e lógica).
\end{itemize}
\end{resumo}

\section*{Exercícios de fixação}
\addcontentsline{toc}{section}{Exercícios de fixação}
\begin{enumerate}[label=\textbf{1.\arabic*}]
  \item Classifique como PL, PLIM, PLI ou PB: (a) escolher quais de 10 projetos financiar; (b) decidir quantas
  unidades inteiras produzir de cada brinquedo e se compra ou não cada molde; (c) escolher valores contínuos
  $x_1,x_2$ e uma binária $y$ que diz qual de duas restrições vale.
  \item Para cada grandeza, diga se uma variável contínua é uma boa aproximação ou se precisamos de inteiros ou
  binários: (a) toneladas de café transportadas; (b) número de caminhões; (c) abrir ou não um armazém;
  (d) litros de combustível.
  \item No exemplo das vans, mude o preço da van para R\$ 300. O ótimo inteiro muda? Arredondar a PL passa a dar o
  ótimo?
\end{enumerate}

\chapter{Por que não basta arredondar}
\label{cap:arredondar}

A primeira ideia de quase todo mundo é: \emph{``resolvo como PL, que é rápido, e arredondo o resultado''}. A
promessa é sedutora: aproveitar a eficiência do simplex e a ``semelhança'' entre o problema inteiro e o contínuo.
Este capítulo mostra, com três exemplos, por que essa ideia falha.

\section{O exemplo de Wolsey}
\label{sec:wolsey}

\[
\max\; x_1 + 0{,}64\,x_2 \quad \sa 50x_1 + 31x_2 \le 250,\quad 3x_1 - 2x_2 \ge -4,\quad x_1,x_2 \ge 0 \text{ inteiros.}
\]

\paragraph{Passo 1 -- resolver a PL (ignorando a inteireza).} O ótimo da PL está num vértice. O vértice
``interessante'' é o encontro das duas retas:
\[
\begin{cases} 50x_1 + 31x_2 = 250\\ 3x_1 - 2x_2 = -4 \end{cases}
\;\Rightarrow\; x_2=\tfrac{3x_1+4}{2}
\;\Rightarrow\; 50x_1 + \tfrac{31(3x_1+4)}{2}=250
\;\Rightarrow\; 193x_1 = 376.
\]
Logo $x_1 = 376/193 \approx 1{,}95$ e $x_2 = 950/193 \approx 4{,}92$, com valor $984/193\approx 5{,}098$. Os outros
vértices, $(0,0)$, $(5,0)$ e $(0,2)$, valem $0$, $5$ e $1{,}28$. Então o ótimo da PL é $(1{,}95;\,4{,}92)$.

\paragraph{Passo 2 -- arredondar.} Os quatro vizinhos inteiros do ótimo da PL são:
\begin{center}
\begin{tabular}{ccccc}
\toprule
ponto & $50x_1+31x_2\le250$? & $3x_1-2x_2\ge-4$? & viável? & valor\\
\midrule
$(1,4)$ & $174$ \checkmark & $-5$ $\times$ & não & --\\
$(1,5)$ & $205$ \checkmark & $-7$ $\times$ & não & --\\
$(2,4)$ & $224$ \checkmark & $-2$ \checkmark & \textbf{sim} & $4{,}56$\\
$(2,5)$ & $255$ $\times$ & $-4$ \checkmark & não & --\\
\bottomrule
\end{tabular}
\end{center}
O arredondamento ``natural'' $(2,5)$ é \textbf{inviável}. O único vizinho viável vale $4{,}56$.

\paragraph{Passo 3 -- o ótimo inteiro de verdade.} Testando todos os pontos inteiros viáveis (são poucos; veja a
figura), o melhor é $(5,0)$, com valor $5$, \textbf{do outro lado} do polígono. O segundo melhor é $(3,3)$, com
$4{,}92$.

\begin{figure}[H]
\centering
\begin{tikzpicture}[scale=0.95]
  \fill[azul!15] (0,0) -- (5,0) -- (1.948,4.922) -- (0,2) -- cycle;
  \draw[->] (-0.3,0) -- (6.3,0) node[right] {$x_1$};
  \draw[->] (0,-0.3) -- (0,6.4) node[above] {$x_2$};
  \foreach \i in {1,...,6} \draw (\i,0.08) -- (\i,-0.08) node[below] {\small\i};
  \foreach \i in {1,...,6} \draw (0.08,\i) -- (-0.08,\i) node[left] {\small\i};
  \draw[thick] (1.16,6.2) node[above left] {\small $50x_1+31x_2=250$} -- (5.2,-0.32);
  \draw[thick] (0,2) -- (2.8,6.2) node[above right] {\small $3x_1-2x_2=-4$};
  \foreach \p in {(0,0),(0,1),(0,2),(1,0),(1,1),(1,2),(1,3),(2,0),(2,1),(2,2),(2,3),(2,4),(3,0),(3,1),(3,2),(3,3),(4,0),(4,1),(5,0)}
    \fill[azul] \p circle (2.2pt);
  \node[star,star points=5,fill=red,inner sep=1.6pt] at (1.948,4.922) {};
  \node[right,red] at (2.4,4.45) {\small ótimo da PL $\approx(1{,}95;\,4{,}92)$};
  \draw[very thick] (1.88,4.88) -- (2.12,5.12) (1.88,5.12) -- (2.12,4.88);
  \node[right] at (2.3,5.2) {\small $(2,5)$ inviável};
  \node[star,star points=5,fill=verde,inner sep=2pt] at (5,0) {};
  \node[right,verde] at (5.2,0.55) {\small ótimo inteiro $(5,0)$};
\end{tikzpicture}
\caption{Exemplo de Wolsey: região viável da PL (azul claro), pontos inteiros viáveis (pontos azuis), ótimo da PL
(estrela vermelha) e ótimo inteiro (estrela verde). O ótimo inteiro está longe do ótimo da PL.}
\label{fig:wolsey}
\end{figure}

\begin{resumo}[Lição]
Arredondar a solução da PL pode dar (1) uma solução \textbf{inviável} ou (2) uma solução viável mas
\textbf{longe do ótimo}. O ótimo inteiro pode estar ``do outro lado'' da região viável.
\end{resumo}

\section{Outro exemplo clássico: o livro do MIT}

Bradley, Hax e Magnanti~\cite[Seç.~9.4]{amp} usam um exemplo igualmente pequeno para mostrar o mesmo fenômeno:
\[
\max\; 5x_1+8x_2\quad\sa x_1+x_2\le6,\quad 5x_1+9x_2\le45,\quad x_1,x_2\ge0\text{ inteiros.}
\]
\begin{center}
\begin{minipage}{0.42\linewidth}
\centering
\begin{tikzpicture}[scale=0.62]
  \fill[verde!15] (0,0) -- (6,0) -- (2.25,3.75) -- (0,5) -- cycle;
  \draw[->] (-0.3,0) -- (7,0) node[right] {$x_1$};
  \draw[->] (0,-0.3) -- (0,6.4) node[above] {$x_2$};
  \foreach \i in {1,...,6} \draw (\i,0.08) -- (\i,-0.08) node[below] {\scriptsize\i};
  \foreach \i in {1,...,6} \draw (0.08,\i) -- (-0.08,\i) node[left] {\scriptsize\i};
  \draw (0,6) -- (6.3,-0.3);
  \draw (-0.2,5.11) -- (7,1.11);
  \foreach \p in {(0,0),(0,1),(0,2),(0,3),(0,4),(0,5),(1,0),(1,1),(1,2),(1,3),(1,4),(2,0),(2,1),(2,2),(2,3),(3,0),(3,1),(3,2),(3,3),(4,0),(4,1),(4,2),(5,0),(5,1),(6,0)}
    \fill[verde!50!black] \p circle (2pt);
  \node[star,star points=5,fill=red,inner sep=1.4pt] at (2.25,3.75) {};
  \draw[very thick] (1.85,3.85) -- (2.15,4.15) (1.85,4.15) -- (2.15,3.85);
  \draw[thick,azul] (2,3) circle (5pt);
  \node[star,star points=5,fill=azul,inner sep=1.8pt] at (0,5) {};
\end{tikzpicture}
\end{minipage}%
\begin{minipage}{0.56\linewidth}
\footnotesize
\begin{tabular}{lccc}
\toprule
& $x_1$ & $x_2$ & $z$\\
\midrule
ótimo da PL ({\color{red}$\bigstar$}) & 2,25 & 3,75 & 41,25\\
arredondado ($\times$) & 2 & 4 & inviável\\
viável mais próximo ($\circ$) & 2 & 3 & 34\\
ótimo inteiro ({\color{azul}$\bigstar$}) & 0 & 5 & \textbf{40}\\
\bottomrule
\end{tabular}
\end{minipage}
\end{center}
O ponto inteiro mais próximo do ótimo da PL é inviável ($5\cdot2+9\cdot4=46>45$), e o ponto viável mais próximo
vale 34, longe do ótimo 40. Os autores observam ainda que, mudando a escala dos coeficientes, dá para construir
exemplos em que o ótimo inteiro fica \textbf{tão longe quanto se queira} da solução arredondada, tanto em valor quanto
em distância no plano.

\section{O caso binário é ainda pior}
\label{sec:mochila3}

\[
\max\; 12x_1 + 7x_2 + 6x_3 \quad \sa 3x_1 + 2x_2 + 2x_3 \le 4,\quad x\in\B^3.
\]
Pense em três itens com valores 12, 7 e 6 e pesos 3, 2 e 2, numa mochila de capacidade 4. Uma estratégia
``gulosa'' razoável é pegar primeiro o item com maior \textbf{razão valor/peso}: $12/3=4$, $7/2=3{,}5$ e $6/2=3$. A
gulosa pega o item 1 (peso 3) e não cabe mais nada: valor \textbf{12}. Mas olhe a enumeração completa:
\begin{center}
\begin{tabular}{ccc|cc|c}
\toprule
$x_1$ & $x_2$ & $x_3$ & peso & valor & viável?\\
\midrule
0&0&0&0&0&sim\\ 0&0&1&2&6&sim\\ 0&1&0&2&7&sim\\ 0&1&1&4&\textbf{13}&\textbf{sim}\\
1&0&0&3&12&sim\\ 1&0&1&5&18&não\\ 1&1&0&5&19&não\\ 1&1&1&7&25&não\\
\bottomrule
\end{tabular}
\end{center}
O ótimo é levar os itens 2 e 3 (valor 13). E ``arredondar'' uma variável binária fracionária significa jogá-la
para 0 ou para 1, ou seja, tomar uma \emph{decisão completamente diferente}. Com binárias não existe ``perto''.

\section{Quando a inteireza é natural: escala de enfermeiras}
\label{sec:enfermeiras}

Cada enfermeira trabalha \textbf{5 dias consecutivos} e folga 2. A demanda de enfermeiras nos dias de domingo a
sábado é $d=(17,13,15,19,14,16,11)$. Queremos contratar o mínimo de enfermeiras.

\paragraph{A escolha da variável é o pulo do gato.} Se usássemos ``número de enfermeiras trabalhando no dia $j$'', não
conseguiríamos garantir os 5 dias consecutivos. A variável certa é
\[
x_k = \text{número de enfermeiras que \textbf{começam} a escala no dia } k,\qquad k=0,\dots,6\ (\text{dom},\dots,\text{sáb}).
\]
Quem começa no dia $k$ trabalha nos dias $k, k+1, \dots, k+4$ (contando ``em círculo'': depois de sábado vem
domingo). Então o dia $j$ é coberto por quem começou em $j, j-1, j-2, j-3$ ou $j-4$. Por exemplo, o domingo ($j=0$)
é coberto por quem começou em dom, sáb, sex, qui e qua:
\[
x_0 + x_3 + x_4 + x_5 + x_6 \ge 17.
\]
O modelo completo é
\[
\min \sum_{k=0}^{6} x_k \quad\sa \sum_{k:\;(j-k)\bmod 7\,\le\,4} x_k \ge d_j\quad (j=0,\dots,6),\qquad x_k\in\Z_+.
\]
A condição $(j-k) \bmod 7 \le 4$ diz exatamente ``o dia $j$ está entre os 5 dias de quem começou em $k$''. No notebook
ela aparece como \texttt{(j - k) \% 7 <= 4}.

\paragraph{Resultados.}
\begin{center}
\begin{tabular}{lcc}
\toprule
& solução & total\\
\midrule
PL (contínua) & $(1{,}33;\ 5{,}33;\ 0;\ 7{,}33;\ 0;\ 3{,}33;\ 5)$ & $22{,}33$\\
PL arredondada para cima & $(2,6,0,8,0,4,5)$ & $25$\\
Ótimo inteiro & $(2,6,0,7,0,3,5)$ & $\mathbf{23}$\\
\bottomrule
\end{tabular}
\end{center}
``6,33 enfermeiras'' não faz sentido. Arredondar \emph{cada} variável para cima garante viabilidade (só
aumentamos a cobertura), mas gasta 2 enfermeiras a mais que o necessário. O ótimo inteiro é 23 e, como a PL vale
22,33, sabemos que \emph{nenhuma} escala usa 22 ou menos (a PL é um limite inferior; veja o
Capítulo~\ref{cap:relaxacao}).

\begin{atencao}[a solução da PL não é única]
A PL das enfermeiras tem vários ótimos com valor $22{,}33$. Outro solver pode devolver outra solução fracionária
(e outro arredondamento). O valor $22{,}33$, porém, é sempre o mesmo.
\end{atencao}

\section*{Exercícios de fixação}
\addcontentsline{toc}{section}{Exercícios de fixação}
\begin{enumerate}[label=\textbf{2.\arabic*}]
  \item No exemplo de Wolsey, verifique que $(3,3)$ é viável e calcule seu valor. Por que ele não aparece entre os
  ``vizinhos'' do ótimo da PL?
  \item Na mochila de 3 itens, a capacidade passa a ser 5. Refaça a gulosa e a enumeração. A gulosa acerta agora?
  \item Escreva a restrição de cobertura da quarta-feira ($j=3$) no modelo das enfermeiras.
\end{enumerate}

\chapter{Relaxação linear, limitantes e qualidade de um modelo}
\label{cap:relaxacao}

Este capítulo apresenta o conceito que mais aparece no notebook, mesmo quando não é citado pelo nome: a
\textbf{relaxação linear}. Ela explica por que o valor de $M$ importa, por que uma formulação ``com mais
restrições'' pode ser melhor e o que o solver faz por baixo dos panos.

\section{Relaxação linear}

\begin{definicao}[relaxação linear]
A \textbf{relaxação linear} de um problema inteiro é o mesmo problema sem as exigências de integralidade:
$x\in\Z_+$ vira $x\ge 0$ e $x\in\B$ vira $0\le x\le 1$. Todo o resto (objetivo e restrições) fica igual.
\end{definicao}

Como a relaxação aceita \emph{todas} as soluções do problema inteiro e mais algumas (as fracionárias), vale sempre:
\begin{center}
\fbox{\begin{minipage}{0.86\linewidth}\centering
Em \textbf{maximização}: \quad $z_{\text{PL}} \ge z_{\text{PI}}$ \quad (a relaxação dá um \textbf{limite superior});\\[2pt]
Em \textbf{minimização}: \quad $z_{\text{PL}} \le z_{\text{PI}}$ \quad (a relaxação dá um \textbf{limite inferior}).
\end{minipage}}
\end{center}
Nos exemplos do capítulo anterior: Wolsey (max) tem $5{,}098 \ge 5$; enfermeiras (min) tem $22{,}33\le 23$.

Duas consequências úteis:
\begin{itemize}
  \item Se a solução ótima da relaxação \textbf{já é inteira}, ela é ótima para o problema inteiro.
  \item Se a relaxação é \textbf{inviável}, o problema inteiro também é.
\end{itemize}

No notebook, a função \py{valor_relaxado(prob)} calcula o valor da relaxação linear de um modelo (ela troca
temporariamente todas as variáveis para contínuas, resolve e restaura; veja a Seção~\ref{sec:auxiliares}).

\section{Limitantes e \emph{gap}}

Durante a resolução, o solver mantém dois números (pensando em maximização):
\begin{itemize}
  \item a \textbf{incumbente}: a melhor solução inteira encontrada até agora, cujo valor é um \textbf{limite
  inferior} para o ótimo (o ótimo é pelo menos isso);
  \item um \textbf{limite superior}, que vem das relaxações (o ótimo é no máximo isso).
\end{itemize}
A diferença relativa entre eles é o \textbf{gap}:
\[
\text{gap} = \frac{\text{limite superior} - \text{incumbente}}{\text{incumbente}}.
\]
Quando o gap chega a zero, o solver \textbf{provou} a otimalidade.

\begin{exemplo}[mochila de 10 projetos do Bloco 6]
A heurística gulosa encontra uma carteira de valor 129 (incumbente), e a relaxação linear diz que não dá para passar
de 161,12 (limite superior). O gap inicial é $(161{,}12-129)/129\approx 25\%$. O ótimo (142) está em algum lugar
entre os dois. ``Fechar'' esse intervalo é o trabalho do \emph{branch-and-bound}.
\end{exemplo}

\section{\emph{Branch-and-bound} em uma página}
\label{sec:bb}

Você não precisa implementar o \emph{branch-and-bound} (B\&B) neste notebook, porque o solver CBC faz isso. Mas
entender a ideia explica \textbf{por que modelos com relaxação forte resolvem mais rápido}.

\begin{enumerate}
  \item Resolva a relaxação linear. Se a solução é inteira, acabou.
  \item Se alguma variável $x_j$ tem valor fracionário $v$, \textbf{ramifique} (\emph{branch}): crie dois
  subproblemas, um com $x_j\le\lfloor v\rfloor$ e outro com $x_j\ge\lceil v\rceil$ (para binárias: $x_j=0$ e
  $x_j=1$). Toda solução inteira está em um dos dois.
  \item Em cada subproblema, resolva a relaxação. \textbf{Pode} (\emph{bound}) o nó se: (a) for inviável; (b) a
  relaxação for inteira, caso em que ela vira candidata a incumbente; ou (c) o limite do nó não for melhor que a
  incumbente.
  \item Repita até não sobrar nó aberto.
\end{enumerate}

\begin{exemplo}[B\&B na mochila de 3 itens]
$\max 12x_1+7x_2+6x_3$ s.a.\ $3x_1+2x_2+2x_3\le4$, $x\in\B^3$. A relaxação de uma mochila se resolve pela
gulosa \emph{fracionando} o último item que não cabe inteiro (Capítulo~\ref{cap:bloco6}).

\begin{center}
\begin{tikzpicture}[
  every node/.style={draw,rounded corners,align=center,font=\footnotesize,fill=white,inner sep=3pt},
  level distance=17mm, edge from parent/.style={draw,-{Stealth}},
  level 1/.style={sibling distance=64mm},
  level 2/.style={sibling distance=33mm},
  level 3/.style={sibling distance=26mm}]
\node {raiz: $15{,}5$\\ $x=(1,\frac12,0)$}
  child {node {$15$\\ $x=(\frac23,1,0)$}
    child {node[fill=verde!15] {$13$, inteira\\ $x=(0,1,1)$\\ \textbf{incumbente}} edge from parent node[draw=none,left,fill=none] {$x_1=0$}}
    child {node[fill=black!8] {inviável\\ (peso 5)} edge from parent node[draw=none,right,fill=none] {$x_1=1$}}
    edge from parent node[draw=none,left,fill=none] {$x_2=1$\ }}
  child {node {$15$\\ $x=(1,0,\frac12)$}
    child {node {$14$\\ $x=(\frac23,0,1)$}
      child {node[fill=black!8] {inviável} edge from parent node[draw=none,left,fill=none] {$x_1=1$}}
      child {node[fill=black!8] {$6\le13$\\ podado} edge from parent node[draw=none,right,fill=none] {$x_1=0$}}
      edge from parent node[draw=none,left,fill=none] {$x_3=1$}}
    child {node[fill=black!8] {$12\le13$\\ podado} edge from parent node[draw=none,right,fill=none] {$x_3=0$}}
    edge from parent node[draw=none,right,fill=none] {\ $x_2=0$}};
\end{tikzpicture}
\end{center}
O ramo $x_2=1$ encontra a incumbente 13. O ramo $x_2=0$ tem limite 15 $>13$, então ainda precisa ser explorado;
dentro dele, todos os nós acabam inviáveis ou com valor $\le 13$. Conclusão: \textbf{13 é ótimo}. Se a relaxação
fosse mais ``apertada'' (limites menores), mais nós seriam podados cedo.
\end{exemplo}

\begin{outros}[de Land e Doig aos solvers atuais]
\begin{itemize}
  \item A ideia de limitar com relaxações lineares e dividir o problema vem de Land e Doig~\cite{landdoig1960}. A
  ramificação em dois ramos, $x_j\le\lfloor v\rfloor$ e $x_j\ge\lceil v\rceil$, descrita acima, é a variante de
  Dakin~\cite{dakin1965}, adotada por praticamente todos os solvers.
  \item Bradley, Hax e Magnanti descrevem o B\&B como uma estratégia de ``dividir para conquistar'': particionar a
  região viável em pedaços mais fáceis e, se preciso, particionar os pedaços de novo~\cite[Seç.~9.5]{amp}.
  \item Os \textbf{planos de corte} de Gomory~\cite{gomory1958} acrescentam desigualdades que cortam a solução
  fracionária da relaxação sem cortar nenhum ponto inteiro. Os solvers modernos, como o CBC usado no notebook,
  combinam as duas ideias (\emph{branch-and-cut}).
  \item O progresso é impressionante. Koch et al.\ mediram, de 2001 a 2020, um ganho de cerca de
  \textbf{50$\times$} nos algoritmos de PLIM e \textbf{20$\times$} no hardware: cerca de \textbf{1000$\times$} no
  total, ou seja, o que levava 15 minutos passou a levar menos de 1 segundo~\cite{koch2022}. Bixby faz um
  histórico dessa evolução desde os anos 1990~\cite{bixby2012}.
\end{itemize}
\end{outros}

\section{Formulações: o mesmo problema, modelos diferentes}
\label{sec:formulacoes}

\begin{definicao}[poliedro e formulação]
Um \textbf{poliedro} é um conjunto da forma $P=\{x\in\R^n : Ax\le b\}$. Um poliedro $P$ é uma \textbf{formulação}
para um conjunto de pontos inteiros $X$ se $X = P\cap\Z^n$, isto é, os pontos inteiros de $P$ são exatamente as
soluções do problema.
\end{definicao}

O mesmo conjunto $X$ admite \textbf{infinitas formulações}: basta ``inclinar'' as restrições sem cortar nem
acrescentar pontos inteiros. Todas dão o mesmo ótimo inteiro, mas \textbf{relaxações diferentes}.

\begin{definicao}[formulação mais forte e formulação ideal]
Se $P_1$ e $P_2$ são formulações para $X$ e $P_1\subset P_2$, dizemos que $P_1$ é \textbf{mais forte} (melhor) que
$P_2$: sua relaxação é mais apertada e dá limites melhores. A formulação \textbf{ideal} é a
\textbf{envoltória convexa} $\operatorname{conv}(X)$, o menor poliedro convexo que contém todos os pontos de
$X$. Nela, todos os vértices são inteiros e a relaxação linear já resolve o problema inteiro.
\end{definicao}

\begin{figure}[H]
\centering
\begin{tikzpicture}[scale=0.9]
  \draw[->] (-0.3,0) -- (5,0) node[right] {$x_1$};
  \draw[->] (0,-0.3) -- (0,4) node[above] {$x_2$};
  \fill[laranja!18] (0,0) -- (4.4,0) -- (2.6,3.5) -- (0,2.7) -- cycle;
  \draw[laranja!80!black,thick] (0,0) -- (4.4,0) -- (2.6,3.5) -- (0,2.7) -- cycle;
  \fill[azul!25] (0,0) -- (4,0) -- (3,2) -- (2,3) -- (1,3) -- (0,2) -- cycle;
  \draw[azul,thick] (0,0) -- (4,0) -- (3,2) -- (2,3) -- (1,3) -- (0,2) -- cycle;
  \foreach \p in {(0,0),(1,0),(2,0),(3,0),(4,0),(0,1),(1,1),(2,1),(3,1),(0,2),(1,2),(2,2),(3,2),(1,3),(2,3)}
     \fill \p circle (2pt);
  \node[laranja!80!black,font=\small,align=left] at (6.6,2.9) {formulação $P_2$\\ (mais fraca)};
  \node[azul,font=\small,align=left] at (6.6,1.4) {formulação $P_1 = \operatorname{conv}(X)$\\ (ideal)};
\end{tikzpicture}
\caption{Duas formulações para o mesmo conjunto de pontos inteiros (pontos pretos). Ambas contêm exatamente os
mesmos pontos inteiros, mas a relaxação de $P_2$ inclui muito mais pontos fracionários.}
\end{figure}

\begin{outros}[para ir além em formulações]
O survey de Vielma~\cite{vielma2015} é a referência moderna sobre como obter formulações mais fortes (e menores),
inclusive para disjunções como as do Bloco 2. As notas de aula do curso 15.083J do MIT (Bertsimas e Schulz) dedicam
as primeiras aulas a formulações e formulações ideais~\cite{ocw15083}.
\end{outros}

\begin{atencao}[``menos restrições'' $\neq$ ``modelo melhor'']
Um modelo com mais restrições pode ser \emph{mais forte}: sua relaxação é menor e o B\&B poda mais. Você verá isso
na prática em dois experimentos do notebook: o valor de $M$ (Blocos 1 e 2) e a formulação agregada
$\times$ desagregada (Bloco 3).
\end{atencao}

\section*{Exercícios de fixação}
\addcontentsline{toc}{section}{Exercícios de fixação}
\begin{enumerate}[label=\textbf{3.\arabic*}]
  \item Em um problema de \emph{minimização}, a relaxação deu 40 e a melhor solução inteira conhecida vale 50. O
  ótimo pode ser 38? Pode ser 45? Pode ser 52? Qual o gap?
  \item A relaxação linear de um problema binário devolveu uma solução com todas as variáveis em 0 ou 1. O que você
  pode concluir?
  \item Por que, no exemplo do B\&B, não foi necessário explorar o ramo $x_3=0$ além do primeiro nó?
\end{enumerate}

\chapter{Como formular modelos com variáveis inteiras}
\label{cap:formular}

\section{O ``algoritmo'' de formulação de Wolsey}

Modelar não é adivinhar: há um roteiro que funciona para quase tudo.

\begin{center}
\begin{tikzpicture}[node distance=7mm, every node/.style={font=\small}]
  \node[draw,rounded corners,fill=azul!8,text width=36mm,align=center] (v) {\textbf{1.} (Re)defina as\\ variáveis};
  \node[draw,rounded corners,fill=azul!8,text width=36mm,align=center,right=of v] (r) {\textbf{2.} Escreva restrições cujos pontos viáveis sejam \emph{exatamente} as soluções do problema};
  \node[draw,rounded corners,fill=azul!8,text width=30mm,align=center,right=of r] (o) {\textbf{3.} Escreva a função objetivo};
  \node[draw,diamond,aspect=2,fill=laranja!10,below=9mm of r,align=center,inner sep=1pt] (d) {Dificuldade?};
  \draw[-{Stealth}] (v) -- (r); \draw[-{Stealth}] (r) -- (o);
  \draw[-{Stealth}] (o) |- (d);
  \draw[-{Stealth}] (d) -| node[pos=0.25,below] {sim: repense as variáveis} (v);
\end{tikzpicture}
\end{center}

\begin{itemize}
  \item \textbf{Variáveis primeiro.} A escolha das variáveis define o que você consegue escrever. No exemplo das
  enfermeiras, ``quantas \emph{começam} no dia $k$'' resolveu o problema dos 5 dias consecutivos.
  \item \textbf{Restrições: nem mais nem menos.} Cada solução do problema real tem que satisfazer todas as
  restrições (senão você corta soluções boas) e cada ponto que satisfaz as restrições tem que ser uma solução
  verdadeira (senão o solver ``trapaceia'').
  \item \textbf{Dificuldade?} Se algo não sai linear, quase sempre falta uma variável (em geral, binária).
\end{itemize}

\begin{receita}[perguntas para fazer antes de escrever qualquer código]
\begin{enumerate}
  \item O que eu \emph{decido}? (variáveis e seus tipos: contínua, inteira, binária)
  \item O que eu quero \emph{otimizar}? (objetivo; cuidado com custos que só existem ``se'' algo acontece)
  \item O que \emph{limita} as decisões? (recursos, demandas, regras de negócio, lógica)
  \item Existe alguma regra do tipo ``se'', ``ou'', ``pelo menos'', ``no máximo'', ``só um''? Então precisaremos de
  binárias (Seções~\ref{sec:bigM} e \ref{sec:logica}).
\end{enumerate}
\end{receita}

\section{Variáveis binárias como indicadores de eventos}

Uma variável binária representa uma \textbf{decisão sim/não} (fazer ou não fazer, comprar ou não comprar) ou
\textbf{sinaliza um evento}:
\[
y = \begin{cases}1 & \text{se o evento ocorre,}\\ 0 & \text{caso contrário.}\end{cases}
\]
O desafio é sempre o mesmo: \textbf{amarrar} a binária às outras variáveis usando apenas restrições lineares. A
ferramenta principal para isso é o \emph{big-M}.

\section{A técnica do \emph{big-M}}
\label{sec:bigM}

\begin{exemplo}[custo fixo, a motivação]
Produzir $x$ unidades custa $K(x)=s+cx$ se $x>0$ e $K(0)=0$: há um custo fixo $s$ (preparação de máquina, compra de
molde, taxa inicial do táxi) que só é pago se produzirmos algo.
\begin{center}
\begin{tikzpicture}[scale=0.8]
  \draw[->] (-0.2,0) -- (5,0) node[right] {$x$};
  \draw[->] (0,-0.2) -- (0,3.6) node[above] {$K(x)$};
  \fill (0,0) circle (2.5pt);
  \draw[thick] (0,1.2) -- (4.5,3.2);
  \draw[fill=white] (0,1.2) circle (2.5pt);
  \node[left] at (0,1.2) {$s$};
  \node[right] at (4.5,3.2) {$s+cx$};
\end{tikzpicture}
\end{center}
A função dá um ``salto'' em $x=0$, então não é linear. Solução: criar $y\in\B$ (1 se produzimos) e escrever
\[
K = s\,y + c\,x, \qquad x \le M\,y,
\]
onde $M$ é a maior produção possível. Se $y=0$, a restrição força $x\le 0$ (não produz, não paga). Se $y=1$, ela
vira $x\le M$, que não atrapalha (produz à vontade e paga $s$).
\end{exemplo}

O raciocínio acima é um caso particular de um resultado geral.

\begin{teorema}[Teorema das variáveis binárias, Sierksma]
\label{teo:binarias}
Seja $f:D\to\R$ uma função (do domínio $D$ das variáveis $x$), $\delta\in\B$ e $M\ge\max\{f(x):x\in D\}$. Então,
para todo $x\in D$,
\[
\big(\delta = 0 \;\Rightarrow\; f(x)\le 0\big) \quad\Longleftrightarrow\quad f(x) - M\delta \le 0 .
\]
\end{teorema}
\begin{proof}[Por que vale (por casos)]
Se $\delta=0$, a restrição da direita é $f(x)\le0$, exatamente o que a implicação exige. Se $\delta=1$, a implicação
não exige nada, e a restrição da direita vira $f(x)\le M$, que é sempre verdadeira porque $M$ é pelo menos o maior
valor de $f$. Nos dois casos, os dois lados dizem a mesma coisa.
\end{proof}

Trocando $\delta$ por $1-\delta$ e $f$ por $-f$ obtemos todas as variações de que precisamos (com $M\ge\max f$ e
$m\le\min f$ no domínio):

\begin{center}
\begin{tabular}{ll}
\toprule
Quero dizer\ldots & Restrição linear\\
\midrule
$\delta=0 \Rightarrow f(x)\le 0$ & $f(x)\le M\delta$\\
$\delta=1 \Rightarrow f(x)\le 0$ & $f(x)\le M(1-\delta)$\\
$\delta=0 \Rightarrow f(x)\ge 0$ & $f(x)\ge m\delta$\\
$\delta=1 \Rightarrow f(x)\ge 0$ & $f(x)\ge m(1-\delta)$\\
\midrule
$x>0 \Rightarrow y=1$ \ (com $0\le x\le M$) & $x\le My$\\
$y=1 \Rightarrow x\ge \ell$ \ (com $x\ge0$) & $x\ge \ell\,y$\\
\bottomrule
\end{tabular}
\end{center}

\begin{exemplo}[lendo uma restrição \emph{big-M} pelos casos]
Variáveis $x_1,x_2\in[0,10]$ e a regra ``se $\delta=1$ então $x_1+x_2\le 5$''. Aqui $f=x_1+x_2-5$, cujo maior
valor é $10+10-5=15$. Logo
\[
x_1+x_2-5 \le 15(1-\delta).
\]
\textbf{Teste:} $\delta=1\Rightarrow x_1+x_2\le 5$ \checkmark. \quad $\delta=0\Rightarrow x_1+x_2\le 20$, que
sempre vale no domínio, então a restrição fica ``desligada'' \checkmark.
\end{exemplo}

\section{Como escolher o valor de $M$}
\label{sec:escolherM}

\begin{receita}[o menor $M$ válido]
$M$ deve ser o \textbf{maior valor que $f$ pode assumir} nas situações em que a restrição deve ficar desligada.
Calcule-o com os limites das variáveis (ou resolvendo uma pequena PL, como no Bloco 2). Use o menor valor que
\emph{nunca corta} uma solução válida.
\end{receita}

\begin{atencao}[$M$ pequeno demais $\times$ grande demais]
\begin{itemize}
  \item \textbf{Pequeno demais:} a restrição ``desligada'' ainda corta soluções válidas. O modelo fica
  \textbf{errado} e o ótimo pode piorar. Ex.: se a fábrica consegue produzir 150 bonecas mas você usa $M=100$,
  proibiu produzir mais de 100.
  \item \textbf{Grande demais:} o modelo continua correto, mas (1) a relaxação linear fica \textbf{fraca}: com
  $M=10^6$, $x\le My$ aceita $y=x/10^6\approx 0$, e o custo fixo quase some da relaxação; e (2) há
  \textbf{problemas numéricos}: o solver considera inteiro qualquer valor a menos de uma tolerância (por exemplo,
  $10^{-6}$). Com $M=10^6$ e $y=10^{-6}$ ``igual a zero'', a restrição permite $x=1$ sem pagar o custo fixo!
\end{itemize}
\end{atencao}

\begin{outros}[o que dizem os fabricantes de solvers]
\begin{itemize}
  \item A Gurobi chama o problema numérico acima de \emph{trickle flow} (``fluxo que vaza''). No exemplo da
  documentação, com $y\le1\,000\,000\,x$ ($x$ binária) e a tolerância de integralidade padrão de $10^{-5}$, o valor
  $x=0{,}0000099999$ é aceito como ``zero'' e permite $y=9{,}9999$. A recomendação é usar o \textbf{menor $M$
  válido} e os limites mais apertados possíveis para as variáveis~\cite{gurobitrickle,gurobibigm}.
  \item A documentação do AMPL recomenda que os limites usados no \emph{big-M} sejam tão apertados quanto possível,
  idealmente da ordem de $\pm10^4$, e nunca maiores que o inverso da tolerância de integralidade~\cite{amplnum}.
  \item Bradley, Hax e Magnanti já davam o mesmo conselho em 1977: em geral é melhor usar o menor valor possível
  de $M$ para evitar dificuldades numéricas~\cite[Seç.~9.2]{amp}.
  \item Solvers comerciais (Gurobi, CPLEX) oferecem \textbf{restrições indicadoras}, em que você escreve
  ``$y=1\Rightarrow f(x)\le0$'' diretamente e o solver cuida da lógica sem um $M$ explícito. No notebook usamos
  PuLP/CBC, então escrevemos o \emph{big-M} à mão, e entender o \emph{big-M} é o que permite usar bem essas
  ferramentas.
\end{itemize}
\end{outros}

\section{Lógica com variáveis binárias}
\label{sec:logica}

Uma \textbf{proposição} é uma afirmação que é verdadeira ou falsa (``o investimento 2 é escolhido'', ``a máquina
opera no modo A''). Associamos a cada proposição simples $P_i$ uma binária $\delta_i$, com $\delta_i=1$ se e somente
se $P_i$ é verdadeira. Proposições compostas usam conectivos:

\begin{center}
\begin{tabular}{llcc|ccccc}
\toprule
conectivo & leitura & símbolo & & \multicolumn{4}{c}{tabela-verdade para $(P_1,P_2)=$}\\
& & & & $(1,1)$ & $(1,0)$ & $(0,1)$ & $(0,0)$\\
\midrule
negação & não $P_1$ & $\neg P_1$ & & 0 & 0 & 1 & 1\\
conjunção & $P_1$ e $P_2$ & $P_1\wedge P_2$ & & 1 & 0 & 0 & 0\\
disjunção inclusiva & $P_1$ ou $P_2$ (ou ambos) & $P_1\vee P_2$ & & 1 & 1 & 1 & 0\\
disjunção exclusiva & exatamente um & $P_1\xor P_2$ & & 0 & 1 & 1 & 0\\
implicação & se $P_1$ então $P_2$ & $P_1\Rightarrow P_2$ & & 1 & 0 & 1 & 1\\
equivalência & $P_1$ se e só se $P_2$ & $P_1\Leftrightarrow P_2$ & & 1 & 0 & 0 & 1\\
\bottomrule
\end{tabular}
\end{center}

\begin{atencao}[a implicação só é falsa num caso]
``Se chover, levo guarda-chuva'' só é \emph{mentira} se chover e eu \emph{não} levar. Se não chover, levar ou não
levar está tudo bem. Por isso $P_1\Rightarrow P_2$ é falsa apenas em $(P_1,P_2)=(1,0)$.
\end{atencao}

\begin{receita}[transformar qualquer regra lógica em restrição: ``proibir as linhas falsas'']
\begin{enumerate}
  \item Escreva a tabela-verdade e veja \textbf{em quais linhas a regra é falsa}.
  \item Procure uma desigualdade linear nas $\delta_i$ que seja violada \emph{exatamente} nessas linhas.
  \item Atalho que sempre funciona para uma cláusula ``ou'': $L_1\vee L_2\vee\cdots$ vira
  $\sum(\text{literais})\ge1$, trocando cada literal negado $\neg P_i$ por $(1-\delta_i)$.
\end{enumerate}
\textbf{Exemplo.} $P_1\Rightarrow P_2$ equivale a $\neg P_1\vee P_2$, que vira $(1-\delta_1)+\delta_2\ge1$, ou seja,
$\delta_1-\delta_2\le 0$. Conferindo: a única linha que viola $\delta_1\le\delta_2$ é $(1,0)$, exatamente a linha
falsa \checkmark.
\end{receita}

\begin{center}
\small
\begin{tabularx}{\linewidth}{>{\raggedright\arraybackslash}X l >{\raggedright\arraybackslash}p{3.1cm}}
\toprule
Proposição & Restrição & Linha(s) proibida(s)\\
\midrule
$\neg P_1$ & $\delta_1=0$ & $\delta_1=1$\\
$P_1\wedge P_2$ & $\delta_1+\delta_2=2$ & todas exceto $(1,1)$\\
$P_1\vee P_2$ & $\delta_1+\delta_2\ge1$ & $(0,0)$\\
$P_1\xor P_2$ & $\delta_1+\delta_2=1$ & $(0,0)$ e $(1,1)$\\
$P_1\Rightarrow P_2$ & $\delta_1-\delta_2\le0$ & $(1,0)$\\
$P_1\Leftrightarrow P_2$ & $\delta_1-\delta_2=0$ & $(1,0)$ e $(0,1)$\\
$P_1\wedge\cdots\wedge P_k\Rightarrow P_{k+1}\vee\cdots\vee P_n$ & $\sum_{i\le k}\delta_i-\sum_{i>k}\delta_i\le k-1$ & $(1,\dots,1,0,\dots,0)$\\
pelo menos / exatamente / no máximo $k$ de $n$ & $\sum\delta_i\ge k$ \ / \ $=k$ \ / \ $\le k$ & --\\
\bottomrule
\end{tabularx}
\end{center}

\paragraph{Por que a fórmula geral funciona?} A regra
$P_1\wedge\cdots\wedge P_k\Rightarrow P_{k+1}\vee\cdots\vee P_n$ só é falsa quando \emph{todos} os $k$ primeiros
são verdadeiros e \emph{todos} os demais são falsos. Nessa linha, o lado esquerdo vale $k-0=k$, e a restrição
$\le k-1$ é violada. Em qualquer outra linha, ou algum dos $k$ primeiros é 0, ou algum dos demais é 1, e o lado
esquerdo fica $\le k-1$.

\begin{outros}[o catálogo de restrições lógicas do livro do MIT]
Bradley, Hax e Magnanti~\cite[Seç.~9.2]{amp} organizam as restrições lógicas em ``famílias'' que você vai
reconhecer no notebook (adaptamos a notação: $y_j=1$ quando a restrição $j$ deve valer):
\begin{itemize}
  \item \textbf{Restrições alternativas} (pelo menos uma de $f_1\le b_1$, $f_2\le b_2$):
  $f_j(x)-B_j(1-y_j)\le b_j$ e $y_1+y_2\ge1$. Como basta uma, dá para economizar uma variável: $y_2=1-y_1$. É
  exatamente a disjunção do Bloco 2.
  \item \textbf{Restrições condicionais}: ``se $f_1(x)>b_1$ então $f_2(x)\le b_2$''. Ela só é violada quando
  $f_1>b_1$ \emph{e} $f_2>b_2$, então equivale a ``$f_1\le b_1$ \textbf{ou} $f_2\le b_2$'': de novo uma disjunção.
  \item \textbf{Alternativas $k$ de $p$}: pelo menos $k$ das $p$ restrições devem valer:
  $f_j(x)-B_j(1-y_j)\le b_j$ para todo $j$, e $\sum_jy_j\ge k$.
  \item \textbf{Alternativas compostas}: a região viável é a união de vários ``blocos'' de restrições; usa-se
  \emph{a mesma} binária para todas as restrições de um bloco.
\end{itemize}
\end{outros}

\begin{outros}[indicador de contagem (Unicamp)]
Nas notas de MC558 da Unicamp~\cite{ravelo}, aparece uma receita útil para \emph{saber} se pelo menos $k$ de $n$
condições binárias $x_i$ estão satisfeitas, com uma binária $z$ que vale 1 exatamente nesse caso:
\[
k-n(1-z)\;\le\;\sum_{i=1}^n x_i\;\le\;k-1+n\,z .
\]
Com $z=1$, a da esquerda exige $\sum x_i\ge k$; com $z=0$, a da direita exige $\sum x_i\le k-1$. É a mesma ideia de
``ida e volta'' do Bloco 5, aplicada a uma soma. Para aprofundar a modelagem lógica, os livros de H.~P.~Williams são
referência clássica~\cite{williams2013,williams2009}.
\end{outros}

\section*{Exercícios de fixação}
\addcontentsline{toc}{section}{Exercícios de fixação}
\begin{enumerate}[label=\textbf{4.\arabic*}]
  \item \textbf{Variável semicontínua.} Uma máquina ou fica desligada ($x=0$) ou opera entre 10 e 50 unidades/h.
  Modele com uma binária $y$.
  \item Com $x_1,x_2\in[0,8]$, modele ``se $\delta=0$ então $x_1-x_2\ge 2$''.
  \item Escreva a restrição para ``se $P_1$ então \emph{não} $P_2$'' e confira pela tabela-verdade.
  \item Dadas três restrições $f_i(x)\le 0$ ($i=1,2,3$) com $f_i\le M_i$ no domínio, modele ``pelo menos duas das
  três devem valer''.
  \item Por que usar $M=10^9$ ``para garantir'' é uma má ideia? Dê os dois motivos.
\end{enumerate}

\chapter{PuLP: do modelo no papel ao código}
\label{cap:pulp}

O notebook usa o \textbf{PuLP 2.9.0}, uma biblioteca Python de código aberto do projeto COIN-OR para
\emph{escrever} modelos~\cite{pulp2011}, junto com o solver \textbf{CBC} (\emph{COIN-OR Branch and Cut}), que já vem
instalado com ela e faz o \emph{branch-and-cut} (Capítulo~\ref{cap:relaxacao}). O PuLP só \emph{descreve} o modelo:
quem resolve é o solver, e trocar de solver (GLPK, HiGHS, Gurobi\ldots) exige mudar apenas a chamada de
\py{solve}. A primeira célula do notebook instala a
versão certa (no Colab, pode ser preciso reiniciar a sessão uma vez).

\section{Anatomia de um modelo}

Vamos escrever o exemplo das vans e ônibus do Capítulo~\ref{cap:pl-pi} em PuLP, linha por linha.

\begin{codigo}
import pulp

# 1) o problema: nome e sentido (LpMinimize ou LpMaximize)
prob = pulp.LpProblem("excursao", pulp.LpMinimize)

# 2) as variáveis: nome, limites e tipo (cat)
v = pulp.LpVariable("vans",   lowBound=0, cat="Integer")
o = pulp.LpVariable("onibus", lowBound=0, cat="Integer")

# 3) a função objetivo: a PRIMEIRA expressão (sem <=, >=, ==) somada ao problema
prob += 250*v + 600*o, "custo"

# 4) as restrições: expressões com <=, >= ou ==, cada uma com um nome único
prob += 12*v + 30*o >= 50, "lugares"

# 5) resolver com o CBC (msg=0 silencia o log)
prob.solve(pulp.PULP_CBC_CMD(msg=0))

# 6) ler a solução
print(pulp.LpStatus[prob.status])          # 'Optimal'
print(v.varValue, o.varValue)              # 2.0 1.0
print(pulp.value(prob.objective))          # 1100.0
\end{codigo}

\begin{description}[style=nextline]
  \item[\py{LpProblem(nome, sentido)}] cria o modelo vazio.
  \item[\py{LpVariable(nome, lowBound, upBound, cat)}] cria uma variável. \py{cat} pode ser
  \py{"Continuous"} (padrão), \py{"Integer"} ou \py{"Binary"}; também existem as constantes
  \py{pulp.LpContinuous} e \py{pulp.LpInteger}. Sem \py{lowBound} a variável pode ser negativa!
  \item[\py{prob += expressão, "nome"}] se a expressão não tem comparação, ela vira a \textbf{função objetivo}; se tem
  (\py{<=}, \py{>=}, \py{==}), vira uma \textbf{restrição}.
  \item[\py{pulp.lpSum(...)}] soma eficiente de muitos termos (equivale a $\sum$).
  \item[\py{prob.status}] 1 = ótimo (\texttt{Optimal}); outros: \texttt{Infeasible}, \texttt{Unbounded},
  \texttt{Not Solved}.
  \item[\py{x.varValue}] valor da variável na solução; \py{pulp.value(prob.objective)} é o valor ótimo.
\end{description}

\section{Índices, somatórios e dicionários}

Modelos reais têm muitas variáveis indexadas. O padrão do notebook é criar listas ou dicionários por
\emph{compreensão} e usar \py{lpSum} com geradores.

\begin{codigo}
n = 3
lucro = [20, 15, 10]
# lista de variáveis x[0], x[1], x[2]
x = [pulp.LpVariable(f"x_{i}", lowBound=0, cat="Integer") for i in range(n)]
# dicionário indexado por pares (i, j)
y = {(i, j): pulp.LpVariable(f"y_{i}_{j}", cat="Binary")
     for i in range(n) for j in range(n)}

prob += pulp.lpSum(lucro[i]*x[i] for i in range(n)), "lucro"     # sum_i p_i x_i
for j in range(n):                                               # uma restrição por j
    prob += pulp.lpSum(y[i, j] for i in range(n)) == 1, f"um_por_coluna_{j}"
\end{codigo}

\begin{receita}[do papel para o código]
\begin{tabular}{ll}
$\sum_{i} p_i x_i$ & \py{pulp.lpSum(p[i]*x[i] for i in range(n))}\\
``para todo $j$'' ($\forall j$) & um laço \py{for j in ...:} com um \py{prob +=} dentro\\
$x_{ij}$ & \py{x[i, j]} (dicionário com chave tupla)\\
$\sum_{k:\,\text{condição}} x_k$ & \py{pulp.lpSum(x[k] for k in K if condição)}\\
\end{tabular}
\end{receita}

\section{Armadilhas comuns em PuLP}

\begin{atencao}[o que o solver \emph{não} aceita]
\begin{itemize}
  \item \textbf{Desigualdades estritas e ``diferente'':} não existem \py{<}, \py{>} nem \py{!=} em PL/PI. Use
  \py{>= 1} no lugar de \py{> 0} para inteiros, e reescreva ``$\ne$'' com binárias.
  \item \textbf{Produtos de variáveis:} \py{x*y} com $x$ e $y$ variáveis \emph{não é linear} (o PuLP dá erro).
  Linearize (Bloco 5 mostra como).
  \item \textbf{\py{if} sobre variáveis:} \py{if x.varValue > 0:} dentro do modelo não funciona, porque na hora de
  construir o modelo a variável ainda não tem valor. \py{if} só pode depender de \emph{dados}. Lógica sobre variáveis
  é feita com binárias e restrições.
  \item \textbf{Nomes repetidos:} duas restrições com o mesmo nome geram erro. Use f-strings com o índice:
  \py{f"capacidade_{i}"}.
  \item \textbf{Esquecer o objetivo:} sem objetivo, o notebook avisa ``Falta definir a função objetivo''.
  \item \textbf{Esquecer restrições:} o modelo pode ficar ilimitado (\texttt{Unbounded}), como acontece no Bloco 2
  se você rodar a verificação antes de completar o \py{TODO}.
\end{itemize}
\end{atencao}

\section{As funções auxiliares do notebook}
\label{sec:auxiliares}

A terceira célula do notebook (``Funções auxiliares'') define ferramentas usadas em todas as verificações. Você não
precisa alterá-las, mas precisa saber o que fazem.

\begin{description}[style=nextline]
  \item[\py{resolver(prob, mostrar=True, variaveis=True)}] resolve com o CBC, imprime o status, o valor do objetivo, o
  tempo e as variáveis não nulas, e \textbf{devolve o valor ótimo} (ou \py{None} se não for ótimo ou se o modelo
  estiver incompleto).
  \item[\py{rodar(construtor, *args, **kwargs)}] chama a sua função construtora (por exemplo,
  \py{modelo_fabrica}), pega o \py{prob} que ela devolve (sozinho ou como primeiro item de uma tupla) e chama
  \py{resolver}. É por isso que todas as funções de exercício terminam com \py{return prob, ...}.
  \item[\py{relaxado(prob)} e \py{valor_relaxado(prob)}] \py{relaxado} é um gerenciador de contexto
  (\py{with}) que transforma todas as variáveis em contínuas e as restaura no fim; \py{valor_relaxado} usa-o para
  devolver o valor da \textbf{relaxação linear}. Binárias continuam com limites $[0,1]$.
  \item[\py{conferir(chave, obtido, dica="")}] compara o valor obtido com o gabarito \py{ESPERADO[chave]} (com
  tolerância relativa de $10^{-4}$) e imprime ``correto'', ``não bate'' (com a dica) ou ``ainda sem solução''.
\end{description}

\begin{exemplo}[o padrão de todos os exercícios]
\begin{codigo}
def modelo_exemplo(dados, opcao=False):
    prob = pulp.LpProblem("exemplo", pulp.LpMaximize)
    x = ...                     # variáveis
    prob += ...                 # objetivo
    prob += ...                 # restrições
    if opcao:                   # variações do exercício (regra extra, simetria...)
        prob += ...
    return prob, x              # rodar() usa o primeiro item

obj = rodar(modelo_exemplo, dados, opcao=True)   # constrói, resolve e imprime
conferir("chave_do_gabarito", obj)               # compara com o gabarito
\end{codigo}
\end{exemplo}

\section*{Exercício de fixação}
\addcontentsline{toc}{section}{Exercício de fixação}
\begin{enumerate}[label=\textbf{5.\arabic*}]
  \item Escreva em PuLP o problema de Wolsey (Seção~\ref{sec:wolsey}) com um parâmetro \py{inteiro} que escolhe entre
  variáveis inteiras e contínuas. (É exatamente o primeiro exercício do notebook; a resposta está no
  Capítulo~\ref{cap:bloco0}.)
\end{enumerate}


\part{Resolvendo o notebook passo a passo}
\chapter{Bloco 0 -- Por que programação inteira? (aquecimento)}
\label{cap:bloco0}

Este bloco revisita os exemplos do Capítulo~\ref{cap:arredondar}, agora no código. Há um único exercício para
completar (o modelo de Wolsey); o resto são células para rodar e interpretar.

\section{0.1 -- Exercício: o modelo de Wolsey em PuLP}

\begin{discussao}
\begin{itemize}
  \item \emph{Qual a ideia por trás de ``resolver o PL e arredondar''?} Aproveitar a eficiência do simplex,
  supondo que o ótimo inteiro esteja ``perto'' do ótimo contínuo.
  \item \emph{O ótimo do PL será inteiro?} Não: é $(1{,}95;\,4{,}92)$. \emph{O arredondamento será viável?}
  $(2,5)$ não é. \emph{Será ótimo?} O único vizinho viável, $(2,4)$, vale $4{,}56$, e o ótimo inteiro é $5$ em
  $(5,0)$. Veja a Seção~\ref{sec:wolsey}.
\end{itemize}
\end{discussao}

\paragraph{O que o exercício pede.} A função já cria o problema e as duas variáveis; o tipo das variáveis depende do
parâmetro \py{inteiro}. Faltam três linhas: o objetivo e as duas restrições.

\paragraph{Passo a passo.}
\begin{enumerate}
  \item Objetivo $\max\; x_1 + 0{,}64x_2$: como o problema foi criado com \py{LpMaximize}, basta somar a expressão.
  \item Restrição $50x_1+31x_2\le250$: expressão com \py{<=}.
  \item Restrição $3x_1-2x_2\ge-4$: expressão com \py{>=}. Não é preciso ``passar para $\le$'': o PuLP aceita
  \py{>=} diretamente.
\end{enumerate}

\begin{codigo}
def modelo_wolsey(inteiro=True):
    cat = pulp.LpInteger if inteiro else pulp.LpContinuous
    prob = pulp.LpProblem("wolsey", pulp.LpMaximize)
    x1 = pulp.LpVariable("x1", lowBound=0, cat=cat)
    x2 = pulp.LpVariable("x2", lowBound=0, cat=cat)
    prob += 1.00*x1 + 0.64*x2, "obj"        # função objetivo
    prob += 50*x1 + 31*x2 <= 250, "c1"      # restrição 1
    prob += 3*x1 - 2*x2 >= -4, "c2"         # restrição 2
    return prob, x1, x2
\end{codigo}

\paragraph{Verificação.} A célula roda o modelo duas vezes, com \py{inteiro=False} (PL) e com \py{inteiro=True}
(PI), usando a função \py{rodar}. Os valores esperados são $5{,}0984$ e $5$. A célula seguinte desenha a
Figura~\ref{fig:wolsey} e lista os quatro vizinhos inteiros do ótimo da PL, como na tabela da
Seção~\ref{sec:wolsey}.

\begin{atencao}
\begin{itemize}
  \item Escrever o objetivo depois das restrições funciona, mas ele precisa ser a única expressão
  \emph{sem} comparação somada ao problema.
  \item Se aparecer ``Falta definir a função objetivo'', a linha do objetivo não foi executada. Lembre-se de rodar a
  célula da função \emph{antes} da célula de verificação.
\end{itemize}
\end{atencao}

\section{0.1 (continuação) -- O caso binário}

A célula da mochila de 3 itens não tem \py{TODO}: ela enumera todas as $2^3=8$ soluções com
\py{itertools.product([0, 1], repeat=3)}, que gera todas as tuplas $(x_1,x_2,x_3)$ com valores 0 ou 1. Para cada uma,
calcula peso e valor e marca se é viável. A saída mostra o que vimos na Seção~\ref{sec:mochila3}: a gulosa (razões
$4$, $3{,}5$, $3$) pega o item 1 e fica em 12; o ótimo é 13, com os itens 2 e 3.

\begin{resumo}[O que levar desta célula]
Enumerar funciona para $n=3$ (8 casos), mas $2^n$ explode muito rápido. Gulosas são rápidas, mas sem garantia. É
por isso que precisamos de modelos e solvers.
\end{resumo}

\section{0.2 -- Escala de enfermeiras}

Também não há \py{TODO}, mas vale a pena ler o código com cuidado, porque ele usa padrões que você vai repetir.

\begin{codigo}
d = [17, 13, 15, 19, 14, 16, 11]          # demanda de dom..sáb

def modelo_enfermeiras(d, inteiro=True):
    prob = pulp.LpProblem("enfermeiras", pulp.LpMinimize)
    cat = pulp.LpInteger if inteiro else pulp.LpContinuous
    x = [pulp.LpVariable(f"x_{k}", lowBound=0, cat=cat) for k in range(7)]
    prob += pulp.lpSum(x), "total"
    for j in range(7):
        prob += pulp.lpSum(x[k] for k in range(7) if (j - k) % 7 <= 4) >= d[j], f"dia_{j}"
    return prob, x
\end{codigo}

\begin{itemize}
  \item \py{x[k]} = enfermeiras que \textbf{começam} no dia $k$ (veja a Seção~\ref{sec:enfermeiras} para entender por
  que essa é a variável certa).
  \item \py{pulp.lpSum(x)} soma todas as variáveis da lista: o total contratado.
  \item O laço cria uma restrição por dia $j$. O gerador soma apenas os $x_k$ cujo turno de 5 dias inclui o dia $j$,
  isto é, $(j-k)\bmod 7\in\{0,1,2,3,4\}$. O operador \texttt{\%} do Python devolve sempre um resto entre 0 e 6, mesmo
  quando $j-k$ é negativo, o que faz a ``volta'' de sábado para domingo funcionar sozinha.
\end{itemize}

\paragraph{Saída e leitura.} PL $=22{,}33$, PI $=23$, PL arredondada para cima $=25$. A mensagem central:
\textbf{arredondar para cima garante viabilidade, mas não otimalidade}; e o limite da PL ($22{,}33$) prova que 23 é o
melhor possível, já que o total é inteiro e precisa ser $\ge 22{,}33$.

\section{O mapa do resto da aula}

A partir do Bloco 1, a inteireza deixa de vir da natureza do problema e passa a ser uma \textbf{ferramenta de
modelagem}. Cada bloco ensina uma ``receita'' de linearização, e todas estão reunidas no
Capítulo~\ref{cap:resumo}:

\begin{center}
\begin{tabular}{cl}
\toprule
Bloco & Técnica\\
\midrule
1 & Implicações \emph{se--então} (custo fixo, produção condicionada)\\
2 & Restrições disjuntivas (\emph{big-M}, escalonamento)\\
3 & Relações lógicas (tabela-verdade, carteira de investimentos)\\
4 & Representação de valores discretos\\
5 & ``Ou exclusivo'' e indicadores (teoria de variáveis lógicas)\\
6 & Problema da mochila binária\\
7 & \emph{Bin packing}\\
\bottomrule
\end{tabular}
\end{center}

\chapter{Bloco 1 -- Implicações \emph{se--então}}
\label{cap:bloco1}

A ferramenta deste bloco é a \textbf{restrição de ativação}:
\[
x>0 \;\Rightarrow\; y=1 \qquad\Longleftrightarrow\qquad x\le M\,y \qquad (M\ge\text{maior valor possível de }x).
\]
Ela é o Teorema~\ref{teo:binarias} com $f(x)=x$: ``se $y=0$ então $x\le0$''.

\section{1A -- Custo fixo: fabricar exige pagar a preparação}

\paragraph{O problema.} Uma fábrica produz Carrinho (A), Boneca (B) e Quebra-cabeça (C). Cada produto exige um
\textbf{molde}, cujo custo fixo é pago uma única vez \emph{se} o produto for fabricado, e zero se não for. Queremos
maximizar o lucro líquido.

\begin{center}
\begin{tabular}{lccc|c}
\toprule
& Carrinho (A) & Boneca (B) & Quebra-cabeça (C) & disponível\\
\midrule
lucro por unidade $p_i$ (R\$) & 20 & 15 & 10 &\\
horas de máquina $h_i$ & 3 & 1 & 2 & $H=200$\\
plástico $q_i$ (kg) & 1 & 1 & 2 & $P=150$\\
custo fixo do molde $s_i$ (R\$) & 100 & 400 & 100 &\\
\bottomrule
\end{tabular}
\end{center}

\begin{discussao}
\begin{enumerate}
  \item \emph{Quais são as variáveis? Só $x_i$ basta?} Não. O custo fixo é ``tudo ou nada'': vale $s_i$ para
  qualquer $x_i>0$ e $0$ para $x_i=0$, o que não é linear em $x_i$ (Seção~\ref{sec:bigM}). Precisamos de
  $x_i\in\Z_+$ (unidades) e $y_i\in\B$ (comprar ou não o molde $i$).
  \item \emph{Como expressar ``se $x_i>0$ então $y_i=1$''?} Com $x_i\le M_iy_i$.
  \item \emph{O que impede $y_i=0$ com $x_i>0$?} A própria restrição: $y_i=0$ força $x_i\le0$. \emph{E nada impede
  $y_i=1$ com $x_i=0$, por que não é problema?} Porque $y_i=1$ custa $s_i>0$ no objetivo. Pagar um molde sem
  produzir só diminui o lucro, então o solver, maximizando, nunca escolhe isso.
  \item \emph{Que valor dar a $M_i$?} O maior $x_i$ que os recursos permitem (detalhes abaixo).
\end{enumerate}
\end{discussao}

\subsection*{Construindo o modelo, passo a passo}

\textbf{Passo 1 -- variáveis.} $x_i\in\Z_+$ = unidades do produto $i$; \ $y_i\in\B$ = 1 se o molde de $i$ é
comprado.

\textbf{Passo 2 -- objetivo.} Lucro das unidades menos custos fixos \emph{pagos}:
\[
\max\;\sum_i \big(p_i x_i - s_i y_i\big) = 20x_A+15x_B+10x_C-100y_A-400y_B-100y_C.
\]

\textbf{Passo 3 -- recursos.}
\[
3x_A + x_B + 2x_C \le 200\quad\text{(horas)},\qquad x_A+x_B+2x_C\le150\quad\text{(plástico)}.
\]

\textbf{Passo 4 -- ativação.} $x_i\le M_iy_i$ para $i=A,B,C$.

\textbf{Passo 5 -- escolher $M_i$.} Mesmo usando \emph{todo} o recurso só com o produto $i$, não dá para produzir mais
que $H/h_i$ (horas) nem mais que $P/q_i$ (plástico). Logo $M_i=\min\{H/h_i,\;P/q_i\}$:
\begin{center}
\begin{tabular}{lccc}
\toprule
& $H/h_i$ & $P/q_i$ & $M_i$\\
\midrule
A & $200/3\approx66{,}7$ & $150$ & $66{,}7$\\
B & $200$ & $150$ & $150$\\
C & $100$ & $75$ & $75$\\
\bottomrule
\end{tabular}
\end{center}
\emph{Refinamento opcional:} como $x_i$ é inteiro, poderíamos usar $\lfloor 66{,}7\rfloor=66$ para A, um
$M$ ainda mais apertado.

\subsection*{O código}

O esqueleto já cria variáveis e restrições de recursos. Faltam o objetivo (TODO 1) e a ativação dentro do laço
(TODO 2), em que \py{Mi} já está calculado.

\begin{codigo}
def modelo_fabrica(M=None, com_regra=False):
    n = len(lucro)
    prob = pulp.LpProblem("fabrica", pulp.LpMaximize)
    x = [pulp.LpVariable(f"x_{i}", lowBound=0, cat="Integer") for i in range(n)]
    y = [pulp.LpVariable(f"y_{i}", cat="Binary") for i in range(n)]

    prob += pulp.lpSum(horas[i]*x[i]    for i in range(n)) <= H, "horas"
    prob += pulp.lpSum(plastico[i]*x[i] for i in range(n)) <= P, "plastico"

    # TODO 1: lucro das unidades - custo fixo dos moldes comprados
    prob += pulp.lpSum(lucro[i]*x[i] - custo_fixo[i]*y[i] for i in range(n)), "lucro_liquido"

    # TODO 2: ativação x_i <= M_i * y_i
    for i in range(n):
        Mi = M if M is not None else min(H/horas[i], P/plastico[i])
        prob += x[i] <= Mi*y[i], f"ativa_{i}"

    if com_regra:                      # exercício 1B (próxima seção)
        prob += x[2] >= 20*y[1], "regra_B_implica_C"
    return prob, x, y
\end{codigo}

O parâmetro \py{M} existe para o experimento seguinte: com \py{M=None} usamos o $M_i$ apertado de cada produto; com
um número, o mesmo $M$ para todos.

\subsection*{Resultado e interpretação}

O ótimo é \textbf{1875}, com $x=(25,125,0)$ e $y=(1,1,0)$: produz 25 carrinhos e 125 bonecas e \textbf{não} produz
quebra-cabeças. Conferindo: horas $3\cdot25+125=200$ e plástico $25+125=150$, os dois recursos esgotados. Lucro
$500+1875-100-400=1875$.

\emph{Por que o quebra-cabeça fica de fora?} Ele usa 2 kg de plástico por R\$ 10 de lucro (R\$ 5/kg), enquanto a boneca
dá R\$ 15/kg. Com o plástico como gargalo, não compensa pagar R\$ 100 de molde para produzir C. Fixando cada
combinação de moldes e otimizando a produção, obtemos:
\begin{center}
\begin{tabular}{ccc|c|ccc}
\toprule
$y_A$ & $y_B$ & $y_C$ & lucro ótimo & $x_A$ & $x_B$ & $x_C$\\
\midrule
1 & 1 & 0 & \textbf{1875} & 25 & 125 & 0\\
0 & 1 & 0 & 1850 & 0 & 150 & 0\\
1 & 1 & 1 & 1775 & 25 & 125 & 0\\
0 & 1 & 1 & 1750 & 0 & 150 & 0\\
1 & 0 & 0 & 1220 & 66 & 0 & 0\\
1 & 0 & 1 & 1130 & 66 & 0 & 1\\
0 & 0 & 1 & 650 & 0 & 0 & 75\\
\bottomrule
\end{tabular}
\end{center}
Repare nas linhas com $y_C=1$ e $x_C=0$: o modelo \emph{permite} comprar o molde sem usar, mas isso só tira R\$ 100
do lucro. Por isso o solver nunca escolhe essa opção.

\begin{atencao}[e a restrição $y_i\le x_i$?]
Em sala apareceu também a restrição $y_A\le x_A$ (``se não produz, não compra o molde''). Ela é válida, porque
$x_A$ é inteiro e $y_A=1$ exige $x_A\ge1$, mas aqui é \textbf{desnecessária}: o custo fixo no objetivo já empurra
$y_i$ para 0 quando $x_i=0$. Ela só seria necessária se $y_i=1$ fosse ``vantajoso'' (por exemplo, um bônus por
abrir) ou se outra restrição dependesse de $y_i$.
\end{atencao}

\section{Experimento: o valor de $M$ importa?}

Todo $M$ \textbf{válido} dá o mesmo ótimo inteiro (1875), mas a \textbf{relaxação linear} muda muito:
\begin{center}
\begin{tabular}{rcc}
\toprule
$M$ & relaxação linear & gap até o ótimo inteiro\\
\midrule
$10^6$ & 2374,95 & 499,95\\
1000 & 2322,50 & 447,50\\
200 & 2112,50 & 237,50\\
apertado ($M_i$) & 2004,17 & 129,17\\
\bottomrule
\end{tabular}
\end{center}

\paragraph{Por que isso acontece?} Na relaxação, $y_i$ pode ser fracionário. Para pagar o mínimo de custo fixo, o
solver escolhe o menor $y_i$ permitido: $y_i=x_i/M_i$. Com $M$ apertado, a relaxação produz $x=(25,125,0)$ com
$y_A=25/66{,}7=0{,}375$ e $y_B=125/150\approx0{,}833$; o custo fixo pago é
$100\cdot0{,}375+400\cdot0{,}833\approx371$, e o valor fica $2375-371\approx2004$. Com $M=10^6$, $y_B$ fica em
$0{,}000125$: o custo fixo praticamente some e a relaxação vira o problema \emph{sem} custo fixo (lucro 2375).
Quanto maior o $M$, pior o limite e mais trabalho para o \emph{branch-and-bound}. Além disso, $M$ enorme gera
problemas numéricos (Seção~\ref{sec:escolherM}).

\section{1B -- Produção condicionada}

\paragraph{A regra.} \emph{Se a fábrica produzir Boneca (B), é obrigada a produzir Quebra-cabeça (C) em lote de pelo
menos 20 unidades.}

\begin{discussao}
\begin{itemize}
  \item \emph{Como forçar $x_C\ge20$ apenas quando $y_B=1$?} Com $x_C\ge 20\,y_B$. Se $y_B=1$: $x_C\ge20$. Se $y_B=0$:
  $x_C\ge0$, que já valia.
  \item \emph{Qual a contrapositiva? Precisa de outra restrição?} A contrapositiva é ``se $x_C<20$, então não
  produz B''. Ela já está garantida: se $x_C<20$, a restrição força $y_B=0$, e a ativação $x_B\le M_By_B$ força
  $x_B=0$. Uma única restrição basta.
  \item \emph{A regra muda o ótimo?} Sim, porque o ótimo de 1A tinha $x_B=125$ e $x_C=0$, que agora é proibido.
\end{itemize}
\end{discussao}

\paragraph{Código.} Dentro de \py{if com_regra:} (já mostrado acima):
\begin{codigo}
    if com_regra:
        prob += x[2] >= 20*y[1], "regra_B_implica_C"     # x_C >= 20 * y_B
\end{codigo}
Lembre-se dos índices: \py{0}=Carrinho, \py{1}=Boneca, \py{2}=Quebra-cabeça.

\paragraph{Resultado.} Ótimo \textbf{1375} com $x=(25,85,20)$ e $y=(1,1,1)$. Agora é preciso comprar o molde de C
(R\$ 100), e cada quebra-cabeça usa recursos que iriam para bonecas: horas $75+85+40=200$, plástico $25+85+40=150$.
Lucro: $500+1275+200-600=1375$.

\begin{atencao}[erros típicos do Bloco 1]
\begin{itemize}
  \item Esquecer $-s_iy_i$ no objetivo: o modelo ``esquece'' o custo fixo e devolve 2375.
  \item Usar $M$ pequeno demais (por exemplo, $M=50$ para todos): corta soluções válidas e o ótimo cai.
  \item Escrever $x_C\ge20$ sem multiplicar por $y_B$: força produzir C \emph{sempre}, e não apenas quando há
  boneca.
  \item Trocar índices (\py{y[2]} no lugar de \py{y[1]}): a regra passa a ser ``se C então C''.
\end{itemize}
\end{atencao}

\chapter{Bloco 2 -- Restrições disjuntivas}
\label{cap:bloco2}

Às vezes queremos que \textbf{apenas uma} entre duas restrições valha (``ou isto, ou aquilo''). O conjunto viável
deixa de ser convexo, e nenhuma PL pura consegue representá-lo. Com uma binária $y$:

\begin{receita}[disjunção ``$f(x)\le0$ OU $g(x)\le0$'']
\[
f(x)\le M_1(1-y),\qquad g(x)\le M_2\,y,\qquad y\in\B.
\]
\begin{itemize}
  \item $y=1$: a primeira vira $f(x)\le0$ (ativa) e a segunda vira $g(x)\le M_2$ (desligada).
  \item $y=0$: a primeira vira $f(x)\le M_1$ (desligada) e a segunda vira $g(x)\le0$ (ativa).
\end{itemize}
$M_1$ precisa ser grande o suficiente para não cortar nada quando $g$ está ativa, e vice-versa.
\end{receita}

\section{2A -- Modos de operação}

\paragraph{O problema.} Uma linha de produção opera no \textbf{modo A} (limite de energia: $4x_1+2x_2\le80$)
\textbf{ou} no \textbf{modo B} (limite de insumo: $2x_1+5x_2\le100$); não precisa respeitar os dois ao mesmo tempo.
O contrato exige $x_1\ge10$. Maximizar $2x_1+3x_2$.

\begin{figure}[H]
\centering
\begin{tikzpicture}[scale=0.13]
  \fill[azul!30] (10,0) -- (20,0) -- (10,20) -- cycle;
  \fill[verde!30] (10,0) -- (50,0) -- (10,16) -- cycle;
  \fill[azul!30,opacity=0.7] (10,0) -- (20,0) -- (10,20) -- cycle;
  \draw[->] (-1,0) -- (58,0) node[right] {$x_1$};
  \draw[->] (0,-1) -- (0,26) node[above] {$x_2$};
  \draw[dashed] (10,0) -- (10,24) node[above] {\small $x_1=10$};
  \draw[azul,thick] (10,20) -- (20,0);
  \draw[verde!60!black,thick] (10,16) -- (50,0);
  \foreach \v in {10,20,30,40,50} \draw (\v,0.6) -- (\v,-0.6) node[below] {\small\v};
  \foreach \v in {10,20} \draw (0.6,\v) -- (-0.6,\v) node[left] {\small\v};
  \node[azul,font=\small,align=left] at (30,17) {modo A: $4x_1+2x_2\le80$};
  \node[verde!50!black,font=\small] at (42,8) {modo B: $2x_1+5x_2\le100$};
  \node[star,star points=5,fill=red,inner sep=1.8pt] at (50,0) {};
  \node[red,font=\small] at (50,-4.5) {ótimo $(50,0)$};
\end{tikzpicture}
\caption{Região viável real do Bloco 2A: a \textbf{união} do triângulo do modo A (azul) com o triângulo do modo B
(verde). A união não é convexa. O ótimo, com valor 100, está no modo B.}
\end{figure}

\begin{discussao}
\begin{enumerate}
  \item \emph{A região é convexa? Onde fica o ótimo?} Não é convexa (é a união de dois triângulos). Avaliando os
  vértices: no modo A, $(10,20)$ vale $80$ e $(20,0)$ vale $40$; no modo B, $(50,0)$ vale $100$ e $(10,16)$ vale $68$.
  O ótimo é $(50,0)$, com valor 100.
  \item \emph{Para que valor de $y$ cada restrição está ativa?} Com $f(x)=4x_1+2x_2-80$ e
  $g(x)=2x_1+5x_2-100$: $y=1$ ativa o modo A e $y=0$ ativa o modo B.
  \item \emph{Qual o menor $M_1$ que nunca corta uma solução válida?} Quando $y=0$ (modo B), a restrição
  $f(x)\le M_1$ precisa aceitar todo ponto do modo B. Então $M_1\ge\max\{f(x): x\in D,\ g(x)\le0\}$, onde
  $D=\{x_1\ge10,\ x_2\ge0\}$. Isso é uma PL!
\end{enumerate}
\end{discussao}

\subsection*{Exercício 2A-1: calcular $M_1$ e $M_2$ resolvendo duas PLs}

\[
M_1 = \max\{\,4x_1+2x_2-80 \;:\; 2x_1+5x_2\le100,\ x_1\ge10,\ x_2\ge0\,\},
\]
\[
M_2 = \max\{\,2x_1+5x_2-100 \;:\; 4x_1+2x_2\le80,\ x_1\ge10,\ x_2\ge0\,\}.
\]
\textbf{À mão}, pelos vértices: para $M_1$, os vértices do triângulo B são $(10,0)$, $(50,0)$ e $(10,16)$, com
$f=-40,\ 120,\ -8$, logo $M_1=120$. Para $M_2$, os vértices do triângulo A são $(10,0)$, $(20,0)$ e $(10,20)$, com
$g=-80,\ -60,\ 20$, logo $M_2=20$.

\textbf{No código}, a função recebe as \emph{funções} \py{obj} e \py{outra} (definidas com \py{lambda}) e as aplica às
variáveis do PuLP. Isso funciona porque \py{4*x1 + 2*x2 - 80} com \py{x1, x2} sendo variáveis do PuLP gera uma
expressão linear.

\begin{codigo}
f = lambda x1, x2: 4*x1 + 2*x2 - 80     # modo A:  f(x) <= 0
g = lambda x1, x2: 2*x1 + 5*x2 - 100    # modo B:  g(x) <= 0

def maximizar_sobre(obj, outra):
    """max obj(x)  s.a.  x em D = {x1 >= 10, x2 >= 0}  e  outra(x) <= 0."""
    prob = pulp.LpProblem("M", pulp.LpMaximize)
    x1 = pulp.LpVariable("x1", lowBound=10)    # o domínio D já está nos limites
    x2 = pulp.LpVariable("x2", lowBound=0)
    prob += obj(x1, x2), "obj"                 # TODO: objetivo
    prob += outra(x1, x2) <= 0, "outra_ativa"  # TODO: a outra restrição está ativa
    return resolver(prob, mostrar=False)

M1 = maximizar_sobre(f, g)     # 120.0
M2 = maximizar_sobre(g, f)     # 20.0
\end{codigo}

\begin{atencao}[o domínio importa]
$M_2$ só é tão pequeno (20) porque usamos $x_1\ge10$. Sem essa informação, o ponto $(0,40)$ seria do modo A e daria
$g=100$. Quanto mais você souber sobre as variáveis, menor (e melhor) o $M$.
\end{atencao}

\subsection*{Exercício 2A-2: o modelo disjuntivo}

\begin{codigo}
def modelo_disjuntivo(M1, M2):
    prob = pulp.LpProblem("disjuntivo", pulp.LpMaximize)
    x1 = pulp.LpVariable("x1", lowBound=10)
    x2 = pulp.LpVariable("x2", lowBound=0)
    y  = pulp.LpVariable("y", cat="Binary")     # y = 1 -> modo A ; y = 0 -> modo B
    prob += 2*x1 + 3*x2, "obj"
    prob += f(x1, x2) <= M1*(1 - y), "modoA"    # TODO
    prob += g(x1, x2) <= M2*y,       "modoB"    # TODO
    return prob, x1, x2, y
\end{codigo}

Resultado: \textbf{100}, em $x=(50,0)$ com $y=0$ (modo B). Antes de completar o TODO, a verificação mostra
``Unbounded'': sem restrições além de $x_1\ge10$ e $x_2\ge0$, o objetivo cresce sem limite.

\begin{atencao}[a chamada \py{rodar(modelo_disjuntivo, M1 or 120, M2 or 20)}]
\py{M1 or 120} significa ``use \py{M1} se ele existir; se for \py{None} (exercício anterior incompleto), use 120''.
Assim a verificação funciona mesmo que o 2A-1 não esteja pronto.
\end{atencao}

\subsection*{Leitura do gráfico: o que $M$ faz com a relaxação}

Eliminando $y$ das duas restrições ($y\ge g/M_2$ e $y\le1-f/M_1$), a região da relaxação linear é
\[
\frac{f(x)}{M_1}+\frac{g(x)}{M_2}\le 1,\qquad f(x)\le M_1,\qquad g(x)\le M_2 .
\]
Quanto maiores os $M$, mais ``frouxa'' essa desigualdade e maior a região laranja do gráfico do notebook:
\begin{center}
\begin{tabular}{ccc}
\toprule
$(M_1,M_2)$ & limite da relaxação & ótimo inteiro\\
\midrule
$(120,20)$ & 100 & 100\\
$(120,100)$ & 116 & 100\\
$(500,500)$ & 285,7 & 100\\
$(5000,5000)$ & 2214,3 & 100\\
\bottomrule
\end{tabular}
\end{center}
Com os $M$ mínimos, a relaxação já dá o ótimo exato (100). Com $M=5000$, o limite passa de 2200. E o que acontece com
$M$ \textbf{menor} que o necessário (o \emph{slider} do notebook)? Com $M_1=50$, por exemplo, no modo B a restrição
``desligada'' vira $4x_1+2x_2\le130$, que corta o ponto $(50,0)$. O ``ótimo'' cai para 82,5: \textbf{modelo errado}.

\section{2B -- Escalonamento em uma máquina}

\paragraph{O problema.} Seis tarefas, com tempo de processamento $p_j$ e peso $w_j$, numa única máquina. Duas
tarefas nunca são processadas ao mesmo tempo. Queremos minimizar a soma ponderada dos tempos de término,
$\sum_j w_jC_j$.
\begin{center}
\begin{tabular}{lcccccc}
\toprule
tarefa & J1 & J2 & J3 & J4 & J5 & J6\\
\midrule
$p_j$ & 4 & 3 & 6 & 2 & 5 & 3\\
$w_j$ & 2 & 1 & 3 & 1 & 2 & 1\\
$p_j/w_j$ & 2 & 3 & 2 & 2 & 2,5 & 3\\
\bottomrule
\end{tabular}
\end{center}

\begin{discussao}
\begin{enumerate}
  \item \emph{Variáveis?} Início $s_j\ge0$ (contínua) e, para cada par $j<k$, uma binária de ordem:
  $y_{jk}=1$ se $j$ vem antes de $k$, e $0$ se $k$ vem antes de $j$.
  \item \emph{Onde está a disjunção?} Para cada par: \textbf{ou} $j$ termina antes de $k$ começar
  ($s_j+p_j\le s_k$), \textbf{ou} $k$ termina antes de $j$ começar ($s_k+p_k\le s_j$).
  \item \emph{Que $M$ vale?} Sem tempo ocioso desnecessário, tudo termina até $\sum_j p_j=23$. Então
  $s_j+p_j-s_k\le23$ sempre, e $M=\sum_jp_j$ basta.
  \item \emph{Existe regra simples?} Sim: a \textbf{regra de Smith} (WSPT). Sequenciar por $p_j/w_j$ crescente é
  ótimo quando não há datas de liberação.
\end{enumerate}
\end{discussao}

\subsection*{O modelo}
\[
\begin{aligned}
\min\;& \sum_j w_j\,(s_j+p_j)\\
\sa & s_j+p_j \le s_k + M(1-y_{jk}) &&\forall\, j<k \quad(\text{se }y_{jk}=1\text{: } j\text{ antes de }k)\\
& s_k+p_k \le s_j + M\,y_{jk} &&\forall\, j<k \quad(\text{se }y_{jk}=0\text{: } k\text{ antes de }j)\\
& s_j\ge0,\quad y_{jk}\in\B.
\end{aligned}
\]
Como $C_j=s_j+p_j$, o objetivo é $\sum_jw_jC_j$. A disjunção vale para \textbf{cada par} $j<k$: são
$\binom{6}{2}=15$ binárias e 30 restrições.

\begin{codigo}
def modelo_1maq(p, w, r=None):
    n = len(p)
    M = sum(p) + (max(r) if r else 0)
    prob = pulp.LpProblem("maquina_unica", pulp.LpMinimize)
    s = [pulp.LpVariable(f"s_{j}", lowBound=0) for j in range(n)]
    y = {(j, k): pulp.LpVariable(f"y_{j}_{k}", cat="Binary")
         for j in range(n) for k in range(j+1, n)}           # y=1: j antes de k

    # TODO 1: objetivo  sum_j w_j * C_j,  C_j = s_j + p_j
    prob += pulp.lpSum(w[j]*(s[j] + p[j]) for j in range(n)), "soma_ponderada"
    # TODO 2: as duas restrições disjuntivas para cada par j < k
    for (j, k), yjk in y.items():
        prob += s[j] + p[j] <= s[k] + M*(1 - yjk), f"j{j}_antes_k{k}"
        prob += s[k] + p[k] <= s[j] + M*yjk,       f"k{k}_antes_j{j}"
    # TODO 3 (desafio): datas de liberação
    if r:
        for j in range(n):
            prob += s[j] >= r[j], f"release_{j}"
    return prob, s, y
\end{codigo}

\paragraph{Lendo o laço.} \py{y.items()} percorre os pares \py{((j, k), variável)}. Para cada par escrevemos as duas
restrições: com $y_{jk}=1$ a primeira está ativa e a segunda fica desligada (folga $M$); com $y_{jk}=0$, o contrário.

\paragraph{Resultado.} Ótimo \textbf{127}. Uma sequência ótima devolvida pelo solver:
\begin{center}
\begin{tikzpicture}[x=5.2mm,y=6mm]
  \foreach \ini/\dur/\nome/\cor in {0/6/J3/azul,6/2/J4/verde,8/4/J1/laranja,12/5/J5/roxo,17/3/J6/cinza,20/3/J2/red!70!black}
  { \fill[\cor!70] (\ini,0) rectangle ++(\dur,1); \draw (\ini,0) rectangle ++(\dur,1);
    \node[white,font=\bfseries\small] at (\ini+\dur/2,0.5) {\nome}; }
  \draw[->] (0,0) -- (24.5,0) node[right] {\small $t$};
  \foreach \t in {0,6,8,12,17,20,23} \node[below,font=\footnotesize] at (\t,0) {\t};
\end{tikzpicture}
\end{center}
Conferindo: $3\cdot6+1\cdot8+2\cdot12+2\cdot17+1\cdot20+1\cdot23=18+8+24+34+20+23=127$. A regra de Smith (ordem
J1, J3, J4, J5, J2, J6) também dá 127. Há empates nas razões, logo várias sequências ótimas.

\begin{outros}[por que a regra de Smith funciona: o argumento da troca]
A regra foi publicada por W.~E.~Smith em 1956~\cite{smith1956}. A prova clássica é por \textbf{troca de vizinhos}.
Suponha uma sequência em que $j$ vem imediatamente antes de $k$, começando no instante $t$, e compare com a sequência
que só troca os dois:
\begin{center}
\begin{tikzpicture}[x=6mm,y=5mm]
  \node[left,font=\small] at (0,0.5) {antes:};
  \fill[black!10] (0,0) rectangle (2,1); \node[font=\scriptsize] at (1,0.5) {\ldots};
  \fill[azul!40] (2,0) rectangle (6,1); \node[font=\small] at (4,0.5) {$j$ ($p_j$)};
  \fill[laranja!40] (6,0) rectangle (9,1); \node[font=\small] at (7.5,0.5) {$k$ ($p_k$)};
  \fill[black!10] (9,0) rectangle (11,1); \node[font=\scriptsize] at (10,0.5) {\ldots};
  \node[below,font=\scriptsize] at (2,0) {$t$};
  \node[left,font=\small] at (0,-1.5) {depois:};
  \fill[black!10] (0,-2) rectangle (2,-1); \node[font=\scriptsize] at (1,-1.5) {\ldots};
  \fill[laranja!40] (2,-2) rectangle (5,-1); \node[font=\small] at (3.5,-1.5) {$k$};
  \fill[azul!40] (5,-2) rectangle (9,-1); \node[font=\small] at (7,-1.5) {$j$};
  \fill[black!10] (9,-2) rectangle (11,-1); \node[font=\scriptsize] at (10,-1.5) {\ldots};
  \node[below,font=\scriptsize] at (2,-2) {$t$};
  \node[right,font=\small,align=left] at (11.5,-0.5) {o resto da sequência\\ não muda};
\end{tikzpicture}
\end{center}
As outras tarefas terminam nos mesmos instantes. Só mudam as contribuições de $j$ e $k$:
\[
\Big(\underbrace{w_j(t+p_j)+w_k(t+p_j+p_k)}_{j\text{ antes de }k}\Big) - \Big(\underbrace{w_k(t+p_k)+w_j(t+p_k+p_j)}_{k\text{ antes de }j}\Big)
= w_kp_j - w_jp_k .
\]
Colocar $j$ antes de $k$ não é pior exatamente quando $w_kp_j\le w_jp_k$, isto é,
$\dfrac{p_j}{w_j}\le\dfrac{p_k}{w_k}$. Numa sequência ótima, então, as razões aparecem em ordem crescente.
O argumento \emph{falha} com datas de liberação: trocar $j$ e $k$ pode obrigar a máquina a esperar.
\end{outros}

\subsection*{Desafio: datas de liberação}

Agora a tarefa $j$ só pode começar em $s_j\ge r_j$, com $r=(5,0,12,0,0,0)$. Basta o TODO 3 acima, mas o $M$ também
precisa crescer: a máquina pode ficar ociosa esperando uma liberação, então usamos $M=\sum p_j+\max r_j$.

\paragraph{Resultado.} Ótimo \textbf{135} (confirmado pela força bruta sobre as $6!=720$ permutações). Uma
sequência ótima:
\begin{center}
\begin{tikzpicture}[x=5.2mm,y=6mm]
  \foreach \ini/\dur/\nome/\cor in {0/2/J4/verde,2/5/J5/roxo,7/4/J1/laranja,11/3/J2/red!70!black,14/6/J3/azul,20/3/J6/cinza}
  { \fill[\cor!70] (\ini,0) rectangle ++(\dur,1); \draw (\ini,0) rectangle ++(\dur,1);
    \node[white,font=\bfseries\small] at (\ini+\dur/2,0.5) {\nome}; }
  \draw[->] (0,0) -- (24.5,0) node[right] {\small $t$};
  \foreach \t in {0,2,7,11,14,20,23} \node[below,font=\footnotesize] at (\t,0) {\t};
  \draw[dashed,red] (5,-0.2) -- (5,1.5) node[above,font=\footnotesize] {$r_1=5$};
  \draw[dashed,red] (12,-0.2) -- (12,1.5) node[above,font=\footnotesize] {$r_3=12$};
\end{tikzpicture}
\end{center}
A regra de Smith deixa de ser ótima: ela quer começar por J1 ou J3, que ainda não foram liberadas. (O valor 127 que
o notebook imprime para ``Smith ignorando $r$'' é o de uma sequência que \emph{viola} as liberações.)

\begin{atencao}[erros típicos do Bloco 2]
\begin{itemize}
  \item Trocar $y$ e $1-y$ entre as duas restrições do mesmo par. Às vezes o modelo ainda funciona por simetria, mas
  a interpretação de $y$ fica invertida, e no 2A o modo errado é ligado.
  \item Usar $M$ que não cobre o maior $r_j$ no desafio: corta sequências válidas.
  \item Escrever a disjunção só para pares consecutivos $(j,j+1)$: ela precisa valer para \textbf{todo} par $j<k$.
  \item No 2A-1, esquecer que $x_1\ge10$ faz parte do domínio: o $M$ calculado fica maior que o necessário.
\end{itemize}
\end{atencao}

\begin{resumo}[Conexão com a lista]
A ideia de ordenar pares de tarefas com binárias e \emph{big-M} é exatamente a das restrições de ordenação do
\emph{flowshop} permutacional (exercício 4 da lista de prática).
\end{resumo}

\chapter{Bloco 3 -- Relações lógicas}
\label{cap:bloco3}

Associamos a cada proposição simples $P_i$ uma binária $\delta_i$ ($\delta_i=1\iff P_i$ verdadeira). Proposições
compostas viram \textbf{restrições lineares sobre os $\delta_i$}. Releia a Seção~\ref{sec:logica} antes de
começar: lá está a técnica de ``proibir as linhas falsas'' da tabela-verdade.

\section{3A -- Laboratório da tabela-verdade}

\paragraph{Como o verificador funciona.} Cada desafio é uma tupla
\py{(nome, nº de variáveis, regra lógica, restrição linear)}. A função \py{verificar_logicas} testa \textbf{todas} as
$2^n$ atribuições de 0/1 e compara: a restrição deve ser satisfeita \emph{exatamente} nas linhas em que a regra é
verdadeira. Você escreve a restrição como uma \py{lambda} que devolve \py{True}/\py{False}; para combinar várias
restrições, use \py{and}.

\begin{discussao}
\begin{itemize}
  \item \emph{$\delta_1+\delta_2\ge1$ representa qual conectivo?} O ``ou inclusivo'': só a linha $(0,0)$ viola.
  \emph{E $\delta_1-\delta_2\le0$?} A implicação $P_1\Rightarrow P_2$: só a linha $(1,0)$ viola.
  \item \emph{Como representar $P_1\Rightarrow P_2$?} Ela só é falsa em $P_1=1, P_2=0$. Precisamos de uma restrição
  violada só nessa linha: $\delta_1\le\delta_2$.
\end{itemize}
\end{discussao}

\subsection*{Derivando cada linha do exercício}

\begin{center}
\small
\begin{tabularx}{\linewidth}{lXl}
\toprule
Regra & Raciocínio (linhas falsas $\to$ restrição) & \py{lambda}\\
\midrule
$\neg P_1$ & falsa em $\delta_1=1$ & \py{a == 0}\\
$P_1\wedge P_2$ & verdadeira só em $(1,1)$: a soma tem que ser 2 & \py{a + b == 2}\\
$P_1\vee P_2$ & falsa só em $(0,0)$: soma $\ge1$ & \py{a + b >= 1}\\
$P_1\xor P_2$ & falsa em $(0,0)$ e $(1,1)$: soma exatamente 1 & \py{a + b == 1}\\
$P_1\Rightarrow P_2$ & $\neg P_1\vee P_2$: $(1-a)+b\ge1$ & \py{a - b <= 0}\\
$P_1\Leftrightarrow P_2$ & falsa em $(1,0)$ e $(0,1)$: iguais & \py{a - b == 0}\\
$P_1\wedge P_2\Rightarrow P_3$ & falsa só em $(1,1,0)$; fórmula geral com $k=2$ & \py{a + b - c <= 1}\\
$P_1\Rightarrow P_2\vee P_3$ & falsa só em $(1,0,0)$; fórmula geral com $k=1$ & \py{a - b - c <= 0}\\
pelo menos 2 de 4 & contar os verdadeiros & \py{a + b + c + d >= 2}\\
exatamente 2 de 4 & idem & \py{a + b + c + d == 2}\\
no máximo 1 de 4 & exclusão mútua & \py{a + b + c + d <= 1}\\
\bottomrule
\end{tabularx}
\end{center}

\begin{exemplo}[conferindo $P_1\wedge P_2\Rightarrow P_3$ pela tabela]
\begin{center}
\begin{tabular}{ccc|c|c|c}
\toprule
$\delta_1$ & $\delta_2$ & $\delta_3$ & regra & $\delta_1+\delta_2-\delta_3$ & $\le1$?\\
\midrule
1&1&1& V & 1 & sim\\
1&1&0& \textbf{F} & \textbf{2} & \textbf{não}\\
1&0&$*$& V & $\le1$ & sim\\
0&1&$*$& V & $\le1$ & sim\\
0&0&$*$& V & $\le0$ & sim\\
\bottomrule
\end{tabular}
\end{center}
A restrição é violada exatamente na única linha falsa \checkmark.
\end{exemplo}

\begin{atencao}[o verificador aceita, mas o solver não]
O verificador avalia expressões Python sobre números 0/1, então aceita coisas como \py{a + b > 0} ou
\py{a - b != 0}, que dão o resultado certo na tabela. Mas \textbf{nenhum solver de PL/PI aceita \py{>} estrito nem
\py{!=}}. Treine escrever a versão que funcionaria no PuLP: \py{a + b >= 1} e \py{a + b == 1}.
\end{atencao}

\section{3B -- Carteira de investimentos}

\paragraph{O problema.} Cinco investimentos ($x_j=1$ se investir em $j$), orçamento 14.
\begin{center}
\begin{tabular}{lccccc}
\toprule
investimento $j$ & 1 & 2 & 3 & 4 & 5\\
\midrule
retorno & 5 & 6 & 4 & 10 & 10\\
custo & 4 & 6 & 5 & 3 & 5\\
\bottomrule
\end{tabular}
\end{center}
Regras: \textbf{R0} orçamento; \textbf{R1} no máximo 3 investimentos; \textbf{R2} investimento 1 \textbf{ou} 2;
\textbf{R3} \textbf{se} 2 \textbf{então} 1; \textbf{R4} \textbf{se} 2, 3 \textbf{ou} 4 \textbf{então} 1; \textbf{R5}
\textbf{se} 1 \textbf{e} 5 \textbf{então} 3.

\subsection*{Traduzindo regra por regra}

\begin{center}
\begin{tabularx}{\linewidth}{lXl}
\toprule
Regra & Lógica & Restrição\\
\midrule
R0 & orçamento & $4x_1+6x_2+5x_3+3x_4+5x_5\le14$\\
R1 & no máximo 3 de 5 & $x_1+x_2+x_3+x_4+x_5\le3$\\
R2 & $P_1\vee P_2$ (inclusivo) & $x_1+x_2\ge1$\\
R3 & $P_2\Rightarrow P_1$ & $x_2\le x_1$\\
R4 & $P_2\Rightarrow P_1$, $P_3\Rightarrow P_1$, $P_4\Rightarrow P_1$ & (a) $x_2+x_3+x_4\le3x_1$ \ ou \ (b) $x_j\le x_1$, $j=2,3,4$\\
R5 & $P_1\wedge P_5\Rightarrow P_3$ ($k=2$) & $x_1+x_5-x_3\le1$\\
\bottomrule
\end{tabularx}
\end{center}

\paragraph{Por que (a) funciona para R4?} Se $x_1=0$, então $x_2+x_3+x_4\le0$ e nenhum dos três pode ser escolhido.
Se $x_1=1$, então $x_2+x_3+x_4\le3$, que sempre vale. O coeficiente 3 é o ``$M$'' (o maior valor possível da
soma).

\begin{codigo}
def modelo_investimentos(desagregada=True):
    prob = pulp.LpProblem("investimentos", pulp.LpMaximize)
    x = {j: pulp.LpVariable(f"x{j}", cat="Binary") for j in J}
    prob += pulp.lpSum(retorno[j]*x[j] for j in J), "retorno"
    prob += pulp.lpSum(custo[j]*x[j] for j in J) <= ORC, "R0_orcamento"
    prob += pulp.lpSum(x[j] for j in J) <= 3, "R1_max3"
    prob += x[1] + x[2] >= 1, "R2_1_ou_2"
    prob += x[2] <= x[1], "R3_2_implica_1"
    if desagregada:                                     # R4 (b): uma por investimento
        for j in (2, 3, 4):
            prob += x[j] <= x[1], f"R4_{j}_implica_1"
    else:                                               # R4 (a): agregada
        prob += x[2] + x[3] + x[4] <= 3*x[1], "R4_agregada"
    prob += x[1] + x[5] - x[3] <= 1, "R5_1e5_implica_3"
    return prob, x
\end{codigo}
Note que o dicionário \py{x} é indexado de 1 a 5 (\py{J = range(1, 6)}), igual ao enunciado.

\paragraph{Resultado.} Ótimo \textbf{21}, com os investimentos $\{1,2,4\}$ (custo $4+6+3=13$). Sem R5, o ótimo seria
25 com $\{1,4,5\}$; R5 exige o 3 junto com 1 e 5, e com o 3 seriam 4 investimentos, o que viola R1.

\begin{atencao}[diagnóstico pelo resultado]
\begin{itemize}
  \item Deu \textbf{26} com $\{2,4,5\}$? Faltam R3/R4: o investimento 2 está sem o 1.
  \item Deu \textbf{25} com $\{1,4,5\}$? Falta (ou está errada) a R5.
  \item Escreveu R5 como $x_1+x_5\le x_3$? Errado: com $x_1=1,\ x_5=0$, isso exige $x_3=1$. A versão certa é
  $x_1+x_5-x_3\le1$.
\end{itemize}
\end{atencao}

\section{Experimento: agregada $\times$ desagregada}

As duas versões de R4 têm \textbf{exatamente as mesmas soluções inteiras}. Qual é ``melhor''?

\begin{center}
\begin{tabular}{lcc}
\toprule
& valor da relaxação & solução da relaxação\\
\midrule
ótimo inteiro & 21 & $\{1,2,4\}$\\
agregada (a) & 23,71 & $x=(0{,}57;\ 0{,}43;\ 0{,}29;\ 1;\ 0{,}71)$\\
desagregada (b) & 22,5 & $x=(0{,}5;\ 0{,}5;\ 0{,}5;\ 0{,}5;\ 1)$\\
\bottomrule
\end{tabular}
\end{center}

\paragraph{Explicação.} Somando as três restrições de (b), $x_2\le x_1$, $x_3\le x_1$ e $x_4\le x_1$, obtemos
$x_2+x_3+x_4\le3x_1$, que é (a). Logo \emph{todo} ponto que satisfaz (b) satisfaz (a): o poliedro de (b) está
\textbf{dentro} do de (a), e (b) é uma formulação \textbf{mais forte} (Seção~\ref{sec:formulacoes}). A recíproca é
falsa: o ponto $x_1=\tfrac13$, $x_2=1$, $x_3=x_4=0$ satisfaz (a) ($1\le1$) mas viola (b) ($1\not\le\tfrac13$).
Com $x_1$ pequeno, (a) ainda deixa um dos outros inteiro; (b) amarra cada um individualmente.

\begin{resumo}[Lições do Bloco 3]
\begin{itemize}
  \item Menor limite superior = relaxação mais forte = menos nós no \emph{branch-and-bound}.
  \item A formulação desagregada tem \emph{mais} restrições e mesmo assim é \emph{melhor}: ``menos restrições''
  não significa ``modelo melhor''.
  \item Com (b), a R3 ($x_2\le x_1$) fica \textbf{redundante}, porque já aparece em R4. Com (a), ela ainda fortalece a
  relaxação.
\end{itemize}
\end{resumo}

\chapter{Bloco 4 -- Representação de valores discretos}
\label{cap:bloco4}

\begin{receita}[variável que só assume valores de uma lista]
Se $x$ só pode valer $v_1,\dots,v_k$ (por exemplo, $\{4,6,8,12,20,24\}$), use uma binária por valor:
\[
x=\sum_{i=1}^{k} v_i\,y_i,\qquad \sum_{i=1}^{k} y_i = 1,\qquad y_i\in\B.
\]
A segunda restrição escolhe \textbf{exatamente um} valor, e a primeira ``lê'' qual foi escolhido.

\textbf{Bônus:} qualquer função de $x$ também fica linear nos $y_i$:
\[
\phi(x)=\sum_{i=1}^{k}\phi(v_i)\,y_i,
\]
\textbf{mesmo que $\phi$ seja não linear}, porque $\phi(v_i)$ é só uma constante para cada $i$.
\end{receita}

\section{O problema: diâmetros de tubulação}

Três trechos de tubulação, com comprimentos $L_j$, precisam transportar uma vazão mínima $q_j$. Cada trecho usa
\textbf{um} diâmetro comercial. A vazão máxima de um tubo de diâmetro $d$ é $0{,}5\,d^2$ L/s (\textbf{não linear} em
$d$). O estoque de cada diâmetro é limitado. Queremos minimizar o custo total.

\begin{center}
\begin{tabular}{lcccccc}
\toprule
diâmetro $d$ (cm) & 4 & 6 & 8 & 12 & 20 & 24\\
\midrule
preço (R\$/m) & 8 & 12 & 18 & 30 & 70 & 100\\
vazão máx.\ $0{,}5d^2$ (L/s) & 8 & 18 & 32 & 72 & 200 & 288\\
estoque (m) & 500 & 500 & 500 & 150 & 250 & 300\\
\bottomrule
\end{tabular}
\qquad
\begin{tabular}{lccc}
\toprule
trecho $j$ & 1 & 2 & 3\\
\midrule
$L_j$ (m) & 300 & 200 & 250\\
$q_j$ (L/s) & 15 & 60 & 150\\
\bottomrule
\end{tabular}
\end{center}

\begin{discussao}
\begin{enumerate}
  \item \emph{Por que não uma variável inteira $d_j$ com $4\le d_j\le24$?} Porque 5, 7, 9, \ldots{} não existem
  comercialmente, e preço e vazão não são funções lineares de $d$.
  \item \emph{A vazão cresce com $d^2$; como o modelo continua linear?} Os coeficientes $0{,}5v_i^2$ são
  \textbf{constantes} (8, 18, 32, \ldots). O que varia são as binárias $y_{jd}$.
  \item \emph{Qual restrição acopla os trechos?} O \textbf{estoque}. Sem ele, cada trecho escolheria sozinho o
  diâmetro mais barato que dá vazão: $(6,12,20)$, custo 27\,100.
\end{enumerate}
\end{discussao}

\section{Construindo o modelo}

\textbf{Variáveis.} $y_{jd}\in\B$ = 1 se o trecho $j$ usa o diâmetro $d$ (3 trechos $\times$ 6 diâmetros = 18
binárias).

\textbf{Objetivo.} O custo de usar $d$ no trecho $j$ é $L_j\cdot\text{preço}_d$ (constante!):
\[
\min\sum_{j}\sum_{d} L_j\,\text{preço}_d\;y_{jd}.
\]

\textbf{Restrições.}
\[
\begin{aligned}
&\textstyle\sum_d y_{jd}=1 &&\forall j &&\text{(um diâmetro por trecho)}\\
&\textstyle\sum_d 0{,}5d^2\,y_{jd}\ge q_j &&\forall j &&\text{(vazão: linear nos }y\text{)}\\
&\textstyle\sum_j L_j\,y_{jd}\le \text{estoque}_d &&\forall d &&\text{(metros usados de cada diâmetro)}
\end{aligned}
\]

\begin{codigo}
def modelo_tubos():
    prob = pulp.LpProblem("tubos", pulp.LpMinimize)
    y = {(j, d): pulp.LpVariable(f"y_{j}_{d}", cat="Binary") for j in range(3) for d in D}
    # TODO 1: custo total
    prob += pulp.lpSum(L[j]*preco[d]*y[j, d] for j in range(3) for d in D), "custo"
    for j in range(3):
        # TODO 2: exatamente um diâmetro por trecho
        prob += pulp.lpSum(y[j, d] for d in D) == 1, f"um_diametro_{j}"
        # TODO 3: vazão mínima (vazao[d] = 0.5*d**2 é um NÚMERO)
        prob += pulp.lpSum(vazao[d]*y[j, d] for d in D) >= q[j], f"vazao_{j}"
    # TODO 4: estoque de cada diâmetro
    for d in D:
        prob += pulp.lpSum(L[j]*y[j, d] for j in range(3)) <= estoque[d], f"estoque_{d}"
    return prob, y
\end{codigo}

\section{Resultado e interpretação}

Para entender a resposta, liste as opções que dão vazão suficiente em cada trecho:
\begin{center}
\begin{tabular}{lll}
\toprule
trecho & diâmetros com vazão suficiente & custo de cada um\\
\midrule
1 ($q=15$, 300 m) & 6, 8, 12, 20, 24 & 3\,600; 5\,400; 9\,000; 21\,000; 30\,000\\
2 ($q=60$, 200 m) & 12, 20, 24 & 6\,000; 14\,000; 20\,000\\
3 ($q=150$, 250 m) & 20, 24 & 17\,500; 25\,000\\
\bottomrule
\end{tabular}
\end{center}
\begin{itemize}
  \item O trecho 2 queria o diâmetro 12, mas só há 150 m em estoque e ele precisa de 200 m. Então o trecho 2 usa 20
  ou 24.
  \item O diâmetro 20 tem 250 m: dá para o trecho 3 (250 m) \emph{ou} para o trecho 2 (200 m), não para os dois.
  \item Opções: trecho 2 com 20 e trecho 3 com 24 custa $14\,000+25\,000=39\,000$; trecho 2 com 24 e trecho 3 com 20
  custa $20\,000+17\,500=37\,500$. A segunda é melhor.
\end{itemize}
Ótimo: diâmetros $(6,\,24,\,20)$, custo $3\,600+20\,000+17\,500=\mathbf{41\,100}$.

\begin{atencao}[erros típicos]
\begin{itemize}
  \item Esquecer o estoque: dá 27\,100 (a dica da verificação avisa).
  \item Usar $\le 1$ em ``um diâmetro por trecho'': o modelo pode não escolher nenhum, mas aí a vazão falha;
  use $=1$, que é o que o problema diz.
  \item Escrever a vazão como \py{0.5*d**2} com \py{d} sendo uma variável do PuLP: não é linear. A potência é
  calculada \emph{sobre os dados}, antes, no dicionário \py{vazao}.
\end{itemize}
\end{atencao}

\begin{resumo}[Para pensar: variável semicontínua]
Uma variável $x\in\{0\}\cup[\ell,u]$ (``ou desligada, ou entre $\ell$ e $u$'') combina este bloco com o Bloco 1:
\[
\ell\,y\le x\le u\,y,\qquad y\in\B.
\]
Com $y=0$, $x=0$; com $y=1$, $\ell\le x\le u$.
\end{resumo}

\chapter{Bloco 5 -- ``Ou exclusivo'' e indicadores (TVL)}
\label{cap:bloco5}

No Bloco 3, as proposições eram \emph{variáveis} ($\delta_1$, $\delta_2$), e $P_1\xor P_2$ virava simplesmente
$\delta_1+\delta_2=1$. Agora as proposições são \textbf{restrições}: queremos ``exatamente uma entre $f(x)\le0$ e
$g(x)\le0$ vale''. Para isso precisamos de binárias que \emph{indiquem} se cada restrição vale, nos dois sentidos.

\section{A receita da teoria de variáveis lógicas (TVL)}

\begin{receita}[indicador $\delta\leftrightarrow f(x)\le0$, com $m\le f(x)\le M$ no domínio]
O ``se e somente se'' tem uma \textbf{ida} e uma \textbf{volta}:
\begin{itemize}
  \item \textbf{Ida} -- $\delta=1\Rightarrow f(x)\le0$:\qquad $f(x)\le M(1-\delta)$.
  \item \textbf{Volta} -- $\delta=0\Rightarrow f(x)>0$. O solver não trabalha com ``$>$'', então usamos
  $f(x)\ge\varepsilon$, onde $\varepsilon>0$ é a menor folga que importa:
  \[
  \varepsilon - f(x) - (\varepsilon - m)\,\delta \le 0 .
  \]
\end{itemize}
Para o ``ou exclusivo'' entre duas restrições: crie $\delta_1\leftrightarrow f\le0$ e $\delta_2\leftrightarrow g\le0$
(ida + volta de cada uma) e acrescente $\delta_1+\delta_2=1$.
\end{receita}

\paragraph{De onde vem a volta?} Queremos ``$\delta=0\Rightarrow \varepsilon-f(x)\le0$''. É o
Teorema~\ref{teo:binarias} aplicado à função $h(x)=\varepsilon-f(x)$: a restrição é $h(x)-M_h\delta\le0$, com
$M_h\ge\max h=\varepsilon-\min f=\varepsilon-m$. Substituindo, chegamos a
$\varepsilon-f(x)-(\varepsilon-m)\delta\le0$. Conferindo pelos casos:
\begin{itemize}
  \item $\delta=0$: fica $\varepsilon-f(x)\le0$, isto é, $f(x)\ge\varepsilon$ \checkmark;
  \item $\delta=1$: fica $m-f(x)\le0$, isto é, $f(x)\ge m$, que sempre vale \checkmark (desligada).
\end{itemize}

\begin{atencao}[notação dos slides]
Nos slides, a volta aparece como ``$\epsilon - f(x) - m_1\delta_1\le0$'', onde $m_1$ faz o papel de
$(\varepsilon-\min f)$. O notebook e esta apostila escrevem $(\varepsilon-m)$ explicitamente, com $m=\min f$, para
evitar ambiguidade. Cuidado ao misturar as duas notações.
\end{atencao}

\section{5A -- Zona de ressonância}

\paragraph{O problema.} A velocidade $x\in\{0,1,\dots,100\}$ de um motor \textbf{não pode ficar na zona de
ressonância} $40\le x\le60$. Há ainda um limite térmico $x\le55$. Maximizar $x$.

\paragraph{A lógica.} ``Fora da zona'' é $(x\le60)\xor(x\ge40)$: exatamente uma das duas vale.
\begin{center}
\begin{tikzpicture}[x=1.1mm]
  \draw[->] (0,0) -- (106,0) node[right] {$x$};
  \fill[red!25] (40,-0.25) rectangle (60,0.25);
  \node[red!70!black,font=\small] at (50,0.65) {zona proibida};
  \draw[azul,thick,|-|] (0,-0.75) -- (60,-0.75) node[midway,below,font=\small] {$x\le60$ ($\delta_1=1$)};
  \draw[verde!60!black,thick,|-|] (40,-1.55) -- (100,-1.55) node[midway,below,font=\small] {$x\ge40$ ($\delta_2=1$)};
  \draw[dashed,laranja] (55,-0.4) -- (55,1.1) node[above,font=\small] {$x\le55$};
  \foreach \t in {0,40,60,100} \node[above,font=\footnotesize] at (\t,0.1) {\t};
\end{tikzpicture}
\end{center}
Para $x<40$, só a primeira vale; para $x>60$, só a segunda; entre 40 e 60, as duas valem, e o ``exclusivo'' proíbe.

\begin{discussao}
\begin{itemize}
  \item \emph{Qual o ótimo?} Com $x\le55$, não dá para pular para $x\ge61$. Fora da zona e abaixo de 55, o melhor é
  $x=39$.
  \item \emph{As regiões $x\le60$ e $x\ge40$ se sobrepõem. Usando só a ``ida'', o modelo proibiria $x=55$?} Não, como
  o experimento vai mostrar.
\end{itemize}
\end{discussao}

\subsection*{Exercício 5A-1: a função \texttt{indicador}}

A função recebe o valor (ou a expressão) de $f$, o indicador $d$, os limites $M$ e $m$ e o $\varepsilon$, e devolve a
\textbf{lista} com as duas restrições:
\begin{codigo}
def indicador(f, d, M, m, eps):
    """Restrições que impõem  d = 1  <=>  f <= 0."""
    ida   = f <= M*(1 - d)                 # d = 1  =>  f <= 0
    volta = eps - f - (eps - m)*d <= 0     # d = 0  =>  f >= eps
    return [ida, volta]
\end{codigo}
O detalhe elegante é que a \textbf{mesma função} serve para números e para o PuLP. O verificador chama
\py{indicador} com números ($f(x)$ calculado, $d\in\{0,1\}$), e a comparação \py{<=} devolve \py{True}/\py{False}.
No modelo, \py{f} e \py{d} são expressões/variáveis do PuLP, e o \py{<=} devolve uma restrição. O verificador testa
todos os $x\in\{0,\dots,100\}$ e $d\in\{0,1\}$ (202 pares) e exige: as restrições valem $\iff$ ($d=1\iff f(x)\le0$).

\subsection*{Exercício 5A-2: o modelo}

\textbf{Limites.} Para $x\in[0,100]$: $f_1=x-60\in[-60,\,40]$ e $f_2=40-x\in[-60,\,40]$. Logo
$M_1=M_2=40$ e $m_1=m_2=-60$. Como $x$ é inteiro, ``$f>0$'' equivale a ``$f\ge1$'': $\varepsilon=1$.

\begin{codigo}
EPS = 1
def modelo_ressonancia(com_volta=True):
    prob = pulp.LpProblem("ressonancia", pulp.LpMaximize)
    x  = pulp.LpVariable("x", lowBound=0, upBound=100, cat="Integer")
    d1 = pulp.LpVariable("d1", cat="Binary")      # d1 = 1 <=> x <= 60
    d2 = pulp.LpVariable("d2", cat="Binary")      # d2 = 1 <=> x >= 40
    prob += 1*x, "velocidade"
    prob += x <= 55, "limite_termico"
    f1 = x - 60
    f2 = 40 - x
    M1, m1 = 40, -60                              # TODO 1
    M2, m2 = 40, -60
    for delta, f_, M_, m_ in ((d1, f1, M1, m1), (d2, f2, M2, m2)):   # TODO 2
        restr = indicador(f_, delta, M_, m_, EPS)
        for c in (restr if com_volta else restr[:1]):   # sem volta: só a ida
            prob += c
    prob += d1 + d2 == 1, "xor"                   # TODO 3
    return prob, x, d1, d2
\end{codigo}

\paragraph{O que essas restrições dizem, em português.} Substituindo os números:
\begin{center}
\begin{tabular}{lll}
\toprule
& ida & volta\\
\midrule
$\delta_1$ & $x-60\le40(1-\delta_1)$: \ se $\delta_1=1$, $x\le60$ & $x\ge61-61\delta_1$: \ se $\delta_1=0$, $x\ge61$\\
$\delta_2$ & $x\ge40\,\delta_2$: \ se $\delta_2=1$, $x\ge40$ & $x\le39+61\delta_2$: \ se $\delta_2=0$, $x\le39$\\
\bottomrule
\end{tabular}
\end{center}
Com $\delta_1+\delta_2=1$ há dois casos. Se $(\delta_1,\delta_2)=(1,0)$, então $x\le60$ e $x\le39$, ou seja,
$x\le39$. Se $(\delta_1,\delta_2)=(0,1)$, então $x\ge61$ e $x\ge40$, ou seja, $x\ge61$, o que é incompatível com
$x\le55$. O ótimo é \textbf{39}.

\section{Experimento: sem a volta}

Só com as idas, o solver devolve \textbf{55}, \emph{dentro} da zona proibida. Ele escolhe $\delta_1=1$ ($x\le60$
\checkmark) e $\delta_2=0$. Sem a volta, $\delta_2=0$ \textbf{não tem consequência nenhuma}: o modelo não ``sabe'' que
$\delta_2=0$ deveria obrigar $x<40$. Com $x=55$, as duas proposições são verdadeiras, violando o ``exclusivo'', mas o
modelo não percebe. A volta ($\delta_2=0\Rightarrow x\le39$) fecha essa brecha.

\section{5B (desafio) -- Desconto por faixa, onde a volta \emph{não} é necessária}

\paragraph{O problema.} Um fornecedor cobra R\$ 10/kg se $x\le99$ kg e R\$ 8/kg (sobre \emph{todos} os kg) se
$x\ge100$ kg. Preciso de pelo menos $D$ kg ($x$ inteiro, até $U=300$). Minimizar o custo.

\begin{discussao}
\begin{enumerate}
  \item \emph{Para $D=90$: comprar 90 kg ou 100 kg?} $90\times10=900$ contra $100\times8=800$: compensa comprar
  \textbf{100 kg}. O ponto de indiferença é $10D=800$, ou seja, $D=80$. A partir de 80 kg, vale comprar 100.
  \item \emph{O custo é $10\,x\,\delta_1+8\,x\,\delta_2$, produto de variável por binária. Como linearizar?} Com
  variáveis auxiliares $u_1=x\delta_1$ e $u_2=x\delta_2$, mais $u_1+u_2=x$, $u_1\le U\delta_1$ e $u_2\le U\delta_2$.
  Como exatamente um $\delta$ vale 1, o $u$ correspondente recebe todo o $x$ e o outro é zero.
  \item \emph{Precisamos da volta?} Não; veja abaixo.
\end{enumerate}
\end{discussao}

\paragraph{Indicadores.} $\delta_1\leftrightarrow x\le99$, com $f_1=x-99\in[-99,\,201]$; \ $\delta_2\leftrightarrow
x\ge100$, com $f_2=100-x\in[-200,\,100]$.

\begin{codigo}
def modelo_desconto(D, U=300, Q=100, EPS=1):
    prob = pulp.LpProblem("desconto", pulp.LpMinimize)
    x  = pulp.LpVariable("x", lowBound=0, upBound=U, cat="Integer")
    d1 = pulp.LpVariable("d1", cat="Binary")      # d1 = 1 <=> x <= Q-1 (preço cheio)
    d2 = pulp.LpVariable("d2", cat="Binary")      # d2 = 1 <=> x >= Q   (desconto)
    u1 = pulp.LpVariable("u1", lowBound=0)        # u1 = x*d1
    u2 = pulp.LpVariable("u2", lowBound=0)        # u2 = x*d2
    prob += 10*u1 + 8*u2, "custo"
    prob += x >= D, "demanda"
    prob += d1 + d2 == 1, "xor"                                   # TODO 1
    f1 = x - (Q - EPS); M1, m1 = U - (Q - EPS), -(Q - EPS)        # TODO 2
    f2 = Q - x;         M2, m2 = Q, Q - U
    for c in indicador(f1, d1, M1, m1, EPS) + indicador(f2, d2, M2, m2, EPS):
        prob += c
    prob += u1 + u2 == x, "soma_u"                                # TODO 3
    prob += u1 <= U*d1, "u1_ativa"
    prob += u2 <= U*d2, "u2_ativa"
    return prob, x
\end{codigo}

\begin{center}
\begin{tabular}{lccccccc}
\toprule
$D$ (kg) & 50 & 75 & 78 & 80 & 81 & 90 & 130\\
\midrule
custo ótimo (R\$) & 500 & 750 & 780 & 800 & 800 & 800 & 1040\\
\bottomrule
\end{tabular}
\end{center}
Ou seja, $\text{custo}(D)=\min\{10D\ (\text{se } D\le99),\ 8\max(D,100)\}$.

\paragraph{Por que a volta é dispensável aqui?} As regiões $x\le99$ e $x\ge100$ \textbf{particionam} o domínio:
todo $x$ está em exatamente uma. Com $\delta_1+\delta_2=1$, ``$\delta_2=0$'' é o mesmo que ``$\delta_1=1$'', e a
\emph{ida} de $\delta_1$ ($x\le99$) já faz o papel da volta de $\delta_2$. Não há como ``trapacear'' escolhendo o
preço com desconto para $x<100$. Na ressonância, as regiões se sobrepunham (entre 40 e 60 as duas valem), e por isso
a volta era indispensável. \emph{Teste:} retire a segunda restrição de \py{indicador} em \py{modelo_desconto}; os
custos não mudam.

\section{5C -- O produto de binárias $\delta_3=\delta_1\delta_2$}

O produto de duas binárias é 1 só quando as duas são 1, ou seja, $\delta_3=1\leftrightarrow(\delta_1=1\wedge\delta_2=1)$,
que é o mesmo que $\delta_3=1\leftrightarrow\delta_1+\delta_2=2$. Cada sentido vira uma restrição pelo Teorema:

\begin{itemize}
  \item \textbf{($\Rightarrow$)} $\delta_3=1\Rightarrow 2-\delta_1-\delta_2\le0$. O indicador da implicação é
  $(1-\delta_3)$ e $\max(2-\delta_1-\delta_2)=2$:
  \[
  2-\delta_1-\delta_2-2(1-\delta_3)\le0 \;\Longrightarrow\; -\delta_1-\delta_2+2\delta_3\le0 .
  \]
  \item \textbf{($\Leftarrow$)} pela contrapositiva, $\delta_3=0\Rightarrow\delta_1+\delta_2-1\le0$, com
  $\max(\delta_1+\delta_2-1)=1$:
  \[
  \delta_1+\delta_2-1-\delta_3\le0 \;\Longrightarrow\; \delta_1+\delta_2-\delta_3\le1 .
  \]
\end{itemize}

\begin{center}
\begin{tabular}{ccc|c|cc}
\toprule
$\delta_1$ & $\delta_2$ & $\delta_3$ & $\delta_3=\delta_1\delta_2$? & $-\delta_1-\delta_2+2\delta_3\le0$ & $\delta_1+\delta_2-\delta_3\le1$\\
\midrule
0&0&0& V & \checkmark & \checkmark\\
0&0&1& F & $\times$ & \checkmark\\
0&1&0& V & \checkmark & \checkmark\\
0&1&1& F & $\times$ & \checkmark\\
1&0&0& V & \checkmark & \checkmark\\
1&0&1& F & $\times$ & \checkmark\\
1&1&0& F & \checkmark & $\times$\\
1&1&1& V & \checkmark & \checkmark\\
\bottomrule
\end{tabular}
\end{center}
As duas restrições juntas são satisfeitas exatamente nas linhas verdadeiras \checkmark. No verificador (no notebook,
os nomes usam letras gregas):
\begin{codigo}
("d3 = d1*d2   (e)",  3, lambda a, b, c: c == (a and b),
     lambda a, b, c: (-a - b + 2*c <= 0) and (a + b - c <= 1)),
("d3 = d1 ou d2 (ou)", 3, lambda a, b, c: c == (a or b),
     lambda a, b, c: (c >= a) and (c >= b) and (c <= a + b)),
\end{codigo}

\paragraph{O ``ou'' (desafio).} $\delta_3=\delta_1\vee\delta_2$: se algum é 1, então $\delta_3=1$ ($\delta_3\ge\delta_1$,
$\delta_3\ge\delta_2$); se os dois são 0, então $\delta_3=0$ ($\delta_3\le\delta_1+\delta_2$).

\begin{atencao}[conexão com o Bloco 3]
$-\delta_1-\delta_2+2\delta_3\le0$ é a soma de $\delta_3\le\delta_1$ e $\delta_3\le\delta_2$: é a versão
\emph{agregada}! As duas restrições separadas são equivalentes nos binários e formam uma relaxação mais forte. Ambas
passam no verificador.

\emph{Observação sobre os slides:} no slide ``Outro exemplo'', a primeira linha traz $\delta_1=1\vee\delta_2=1$, mas a
derivação que segue é a da \textbf{conjunção} ($\wedge$), que é a correta para o produto.
\end{atencao}

\begin{receita}[produto de contínua por binária: $u=x\,\delta$, com $0\le x\le U$]
\[
u\le U\delta,\qquad u\le x,\qquad u\ge x-U(1-\delta),\qquad u\ge0.
\]
Se $\delta=0$: $u\le0$, logo $u=0$. Se $\delta=1$: $u\le x$ e $u\ge x$, logo $u=x$. (No 5B bastaram duas
restrições porque $u_1+u_2=x$ e $\delta_1+\delta_2=1$ já fazem o resto.)
\end{receita}

\begin{atencao}[erros típicos do Bloco 5]
\begin{itemize}
  \item Na volta, usar $m$ no lugar de $(\varepsilon-m)$; ou trocar o sinal de $f$. O verificador mostra um
  contraexemplo $(x,\delta)$; refaça os dois casos ($\delta=0$ e $\delta=1$) à mão.
  \item Calcular $M$ e $m$ errados: lembre que $f_2=40-x$ é \emph{decrescente} em $x$. O máximo ocorre em $x=0$ e o
  mínimo em $x=100$.
  \item Com variáveis contínuas, $\varepsilon$ é uma escolha de modelagem (tolerância). Valores muito pequenos causam
  problemas numéricos.
\end{itemize}
\end{atencao}

\chapter{Bloco 6 -- Problema da mochila binária}
\label{cap:bloco6}

Temos itens $j$ com \textbf{valor} $c_j$ e \textbf{peso} $a_j$, e uma mochila de capacidade $b$. Queremos o conjunto
de maior valor total que cabe. A mochila é um dos problemas clássicos mais importantes: aparece sozinha (orçamento
de projetos, carga de contêineres) e como \emph{subproblema} de modelos maiores.

\section{Quatro versões do problema}

\begin{center}
\begin{tabular}{ccccl}
\toprule
Caso & Fracionável? & Limite por item & Tipo & Solução\\
\midrule
1 & sim & ilimitado & PL & trivial: só o item de maior $c_j/a_j$, $x_{j^*}=b/a_{j^*}$\\
2 & não & ilimitado & PI & \textbf{não trivial}\\
3 & sim & $\le1$ & PL & gulosa por $c_j/a_j$, fracionando o último item\\
\textbf{4} & \textbf{não} & $\le\mathbf{1}$ & \textbf{PB} & \textbf{não trivial} $\leftarrow$ \emph{este bloco}\\
\bottomrule
\end{tabular}
\end{center}

\begin{exemplo}[caso 3 nos slides: a gulosa fracionária]
$\max 6x_1+8x_2+4x_3+x_4$ s.a.\ $x_1+2x_2+2x_3+x_4\le4$, $0\le x_j\le1$. Razões: $6$, $4$, $2$, $1$.
Pegamos $x_1=1$ (sobram 3), $x_2=1$ (sobra 1) e $x_3=1/2$ (acabou); $x_4=0$. Valor $6+8+2=16$.
\end{exemplo}

O caso 3 é exatamente a \textbf{relaxação linear} do caso 4. Então a gulosa fracionária dá, sem simplex nenhum, o
limite superior da mochila binária.

\begin{outros}[Dantzig (1957) e o ``item crítico'']
Esse resultado é clássico: Dantzig mostrou em 1957 que a relaxação contínua da mochila é resolvida ordenando os itens
por $c_j/a_j$ e enchendo a mochila até o primeiro item que não cabe inteiro, o \textbf{item crítico}, que entra
fracionado~\cite{dantzig1957,martellototh1990}. O valor resultante (o ``limite de Dantzig'') é o limite superior
usado nos primeiros algoritmos de \emph{branch-and-bound} para a mochila.
\end{outros}

\section{O problema do notebook: carteira de projetos}

Dez projetos, orçamento $B=74$ (R\$ mil):
\begin{center}
\small
\begin{tabular}{l*{10}{c}}
\toprule
projeto $j$ & 1&2&3&4&5&6&7&8&9&10\\
\midrule
valor $c_j$ & 29&47&14&33&37&13&41&36&38&16\\
custo $a_j$ & 34&25&19&32&18&12&16&36&16&16\\
$c_j/a_j$ & 0,85&1,88&0,74&1,03&2,06&1,08&\textbf{2,56}&1,00&2,38&1,00\\
\bottomrule
\end{tabular}
\end{center}

\begin{discussao}
\begin{enumerate}
  \item \emph{Variável? Por que binária e não inteira qualquer?} $x_j=1$ se o projeto $j$ é financiado. Um projeto
  não pode ser financiado ``duas vezes'' nem pela metade: é o caso 4.
  \item \emph{Objetivo e restrição?} $\max\sum_jc_jx_j$ s.a.\ $\sum_ja_jx_j\le74$. Uma única restrição!
  \item \emph{Qual entra primeiro pela razão? A gulosa resolve? E a relaxação, é o caso 3?} O projeto 7 (razão 2,56)
  entra primeiro. A gulosa \emph{não} garante o ótimo. A relaxação é o caso 3.
\end{enumerate}
\end{discussao}

\subsection*{Exercício 6A}
\begin{codigo}
def modelo_mochila(valor, peso, B, extras=False):
    prob = pulp.LpProblem("mochila", pulp.LpMaximize)
    x = {j: pulp.LpVariable(f"x_{j}", cat="Binary") for j in valor}
    prob += pulp.lpSum(valor[j]*x[j] for j in valor), "valor_total"       # TODO 1
    prob += pulp.lpSum(peso[j]*x[j] for j in valor) <= B, "capacidade"    # TODO 2
    if extras:
        ...                                                               # 6B, abaixo
    return prob, x
\end{codigo}
\py{valor} e \py{peso} são dicionários com chaves $1,\dots,10$ (\py{enumerate(..., start=1)}), então
\py{for j in valor} percorre os projetos.

\section{Comparando quatro respostas}

\paragraph{Gulosa (heurística).} Ordem por razão: 7, 9, 5, 2, 6, 4, 8, 10, 1, 3.
\begin{center}
\begin{tabular}{clccl}
\toprule
passo & projeto & custo acumulado & valor & decisão\\
\midrule
1 & 7 & 16 & 41 & entra\\
2 & 9 & 32 & 79 & entra\\
3 & 5 & 50 & 116 & entra\\
4 & 2 & 75 & -- & não cabe ($75>74$)\\
5 & 6 & 62 & 129 & entra\\
6+ & 4, 8, 10, 1, 3 & -- & -- & nenhum cabe nos 12 restantes\\
\bottomrule
\end{tabular}
\end{center}
Resultado: $\{5,6,7,9\}$, valor \textbf{129}.

\paragraph{Relaxação linear (Dantzig).} Mesma ordem, mas o projeto 2 entra \emph{fracionado} para usar os 24 que
sobram: $x_2=24/25=0{,}96$. Valor $116+0{,}96\cdot47=\mathbf{161{,}12}$. O projeto 2 é o item crítico.

\paragraph{Ótimo (PuLP/CBC) e programação dinâmica.} $\{2,7,9,10\}$, custo $25+16+16+16=73$, valor
$47+41+38+16=\mathbf{142}$. A programação dinâmica $O(nB)$ do notebook confirma (você verá PD mais adiante na
disciplina).

\begin{center}
\begin{tikzpicture}[x=0.075cm]
  \draw[->] (120,0) -- (168,0);
  \foreach \v in {120,130,140,150,160} \draw (\v,0.1) -- (\v,-0.1) node[below,font=\footnotesize] {\v};
  \fill[laranja] (129,0) circle (3pt) node[above=3pt,font=\small,align=center] {gulosa\\129};
  \fill[verde] (142,0) circle (3pt) node[above=3pt,font=\small,align=center] {ótimo\\142};
  \fill[azul] (161.12,0) circle (3pt) node[above=3pt,font=\small,align=center] {relaxação\\161,12};
  \draw[decorate,decoration={brace,mirror,amplitude=5pt}] (129,-0.6) -- (161.12,-0.6)
    node[midway,below=6pt,font=\small] {intervalo que o \emph{branch-and-bound} precisa fechar};
\end{tikzpicture}
\end{center}

\section{6B -- Mochila com regras de negócio}

Juntando os blocos 1 a 3:
\begin{center}
\begin{tabular}{lll}
\toprule
Regra & Lógica & Restrição\\
\midrule
(a) projetos 2 e 7 incompatíveis & exclusão mútua & $x_2+x_7\le1$\\
(b) se 9 então 5 & implicação & $x_9\le x_5$\\
(c) no máximo 4 projetos & cardinalidade & $\sum_jx_j\le4$\\
(d) 3 ou 6 (ou ambos) & ou inclusivo & $x_3+x_6\ge1$\\
\bottomrule
\end{tabular}
\end{center}
\begin{codigo}
    if extras:
        prob += x[2] + x[7] <= 1, "excl_2_7"          # (a)
        prob += x[9] <= x[5], "c9_implica_5"           # (b)
        prob += pulp.lpSum(x.values()) <= 4, "max4"   # (c)
        prob += x[3] + x[6] >= 1, "c3_ou_6"            # (d)
\end{codigo}
\paragraph{Resultado.} Ótimo \textbf{135} com $\{2,5,6,9\}$ (custo $25+18+12+16=71$). A regra (a) corta a solução de
6A (que tinha 2 e 7 juntos); (b) obriga o 5 a acompanhar o 9; (d) é satisfeita pelo 6.

\begin{atencao}[erros típicos]
\begin{itemize}
  \item Escrever (b) como $x_5\le x_9$: isso é ``se 5 então 9'', a implicação ao contrário.
  \item Escrever (d) como $x_3+x_6=1$: isso é o ``ou exclusivo'' e proíbe escolher os dois.
\end{itemize}
\end{atencao}

\begin{resumo}[Reflexão do notebook]
A gulosa fica em 129, a relaxação diz ``até 161'' e o ótimo é 142. Quem fecha a diferença é o
\emph{branch-and-bound}, que usa a relaxação como limite e ramifica nas variáveis fracionárias (aqui, primeiro em
$x_2$).
\end{resumo}

\chapter{Bloco 7 -- \emph{Bin packing}}
\label{cap:bloco7}

Empacotar \textbf{todos} os itens no \textbf{menor número de caixas} de capacidade $b$. É a mochila ``ao
contrário'': agora todo item entra, e o que queremos minimizar é o número de caixas.

\paragraph{Os dados.} Dez itens de tamanhos $w=(60,36,53,51,58,48,47,48,53,35)$, caixas de capacidade $b=100$. A soma
dos tamanhos é 489, então o \textbf{limite de volume} é $\lceil489/100\rceil=5$ caixas.

\begin{discussao}
\begin{enumerate}
  \item \emph{Quais variáveis?} $x_{ij}=1$ se o item $j$ vai para a caixa $i$; $y_i=1$ se a caixa $i$ é usada.
  \item \emph{Quantas caixas considerar?} No pior caso, uma por item: $n$ caixas ($i=1,\dots,n$).
  \item \emph{As duas famílias de restrições?} Todo item vai para \textbf{exatamente uma} caixa, e a capacidade é
  respeitada, \textbf{mas só se a caixa for usada}. ``Só se usada'' é o Bloco 1 de novo:
  $\sum_jw_jx_{ij}\le b\,y_i$.
  \item \emph{Dá para usar 5 caixas?} Não; veja o argumento abaixo.
\end{enumerate}
\end{discussao}

\section{O modelo}
\[
\begin{aligned}
\min\;& \sum_{i=1}^{n} y_i\\
\sa & \sum_{i=1}^{n} x_{ij} = 1 &&\forall j &&\text{(todo item em exatamente uma caixa)}\\
& \sum_{j=1}^{n} w_j x_{ij} \le b\,y_i &&\forall i &&\text{(capacidade, só se a caixa é usada)}\\
& x_{ij},\,y_i\in\B.
\end{aligned}
\]
Repare que $b\,y_i$ faz dois papéis ao mesmo tempo: se $y_i=1$, limita a ocupação a $b$; se $y_i=0$, proíbe
colocar qualquer item na caixa $i$ (é uma restrição de ativação com $M=b$).

\begin{codigo}
def modelo_binpacking(w, b, simetria=False):
    n = len(w)
    prob = pulp.LpProblem("binpacking", pulp.LpMinimize)
    y = {i: pulp.LpVariable(f"y_{i}", cat="Binary") for i in range(n)}
    x = {(i, j): pulp.LpVariable(f"x_{i}_{j}", cat="Binary") for i in range(n) for j in range(n)}
    prob += pulp.lpSum(y.values()), "n_caixas"                                   # TODO 1
    for j in range(n):                                                           # TODO 2
        prob += pulp.lpSum(x[i, j] for i in range(n)) == 1, f"item_{j}_em_uma_caixa"
    for i in range(n):                                                           # TODO 3
        prob += pulp.lpSum(w[j]*x[i, j] for j in range(n)) <= b*y[i], f"capacidade_{i}"
    if simetria:                                                                 # desafio
        for i in range(n-1):
            prob += y[i] >= y[i+1], f"simetria_{i}"
    return prob, x, y
\end{codigo}

\section{Resultado: 6 caixas}

\begin{center}
\begin{tikzpicture}[x=1.45cm,y=0.032cm]
  \foreach \cx/\a/\b/\la/\lb in {1/60/35/60/35, 2/58/0/58/, 3/36/53/36/53, 4/51/48/51/48, 5/47/53/47/53, 6/48/0/48/}{
    \draw[thick] (\cx-0.35,0) rectangle (\cx+0.35,100);
    \fill[azul!35] (\cx-0.35,0) rectangle (\cx+0.35,\a);
    \node[font=\small] at (\cx,\a/2) {\la};
    \ifnum\b>0
      \fill[verde!35] (\cx-0.35,\a) rectangle (\cx+0.35,\a+\b);
      \node[font=\small] at (\cx,\a+\b/2) {\lb};
    \fi
    \node[below,font=\footnotesize] at (\cx,0) {caixa \cx};
  }
  \draw[dashed,red] (0.5,100) -- (6.5,100) node[right,font=\footnotesize] {$b=100$};
\end{tikzpicture}
\end{center}
Ocupações: 95, 58, 89, 99, 100 e 48. A heurística FFD (\emph{First-Fit Decreasing}) do notebook também usa 6 caixas,
e a relaxação linear vale $4{,}89$.

\paragraph{Por que é impossível usar 5 caixas?} Um argumento curto, só com aritmética:
\begin{enumerate}
  \item Há \textbf{cinco itens maiores que 50}: 60, 58, 53, 53 e 51. Dois deles nunca cabem juntos, então cada um
  ocupa uma caixa diferente. Com 5 caixas, cada caixa teria exatamente um item grande.
  \item As folgas dessas caixas seriam $40,\ 42,\ 47,\ 47$ e $49$.
  \item Os itens pequenos são 48, 48, 47, 36 e 35. Dois pequenos juntos somam pelo menos $35+36=71>49$, então cada
  caixa recebe \emph{no máximo um} pequeno.
  \item Os \textbf{dois itens de tamanho 48} precisam de folga $\ge48$, mas só a caixa do 51 tem folga 49. Um dos 48
  fica sem lugar.
\end{enumerate}
Logo são necessárias pelo menos 6 caixas, e o solver encontrou 6: \textbf{6 é ótimo}.

\begin{outros}[limitantes para o \emph{bin packing}]
O passo 1 do argumento é a ideia por trás dos limitantes inferiores clássicos de Martello e Toth: itens maiores que
metade da capacidade não podem dividir caixa entre si e precisam ser contados separadamente, o que melhora o limite
de volume $\lceil\sum w_j/b\rceil$~\cite{martellototh1990,delorme2016}. É um exemplo de como raciocínio
combinatório complementa a relaxação linear.
\end{outros}

\section{Reflexões: relaxação fraca e simetria}

\paragraph{Relaxação fraca.} Na relaxação, o solver pode usar $y_i=\sum_jw_jx_{ij}/b$, ``meia caixa'' aqui e ali, e
o valor fica exatamente $\sum_jw_j/b=4{,}89$. A relaxação não ``enxerga'' que itens grandes não se combinam.

\begin{outros}[a formulação ``de Kantorovich'' e alternativas]
O modelo com $x_{ij}$ e $y_i$ é atribuído a Kantorovich~\cite{kantorovich1960}. Os surveys de Valério de
Carvalho~\cite{carvalho2002} e de Delorme, Iori e Martello~\cite{delorme2016} registram que sua relaxação linear é
fraca (vale apenas o limite de volume) e que formulações baseadas em \emph{padrões de corte}, como a de Gilmore e
Gomory~\cite{gilmoregomory1961}, têm relaxações muito mais fortes, ao custo de um número exponencial de variáveis.
\end{outros}

\paragraph{Simetria.} As caixas são \textbf{idênticas}: trocar os rótulos de duas caixas gera uma solução
``diferente'' para o solver, mas igual na prática. Uma solução com $k$ caixas tem $k!$ cópias equivalentes, e o
\emph{branch-and-bound} pode perder tempo explorando todas.

\paragraph{Quebra de simetria.} Exigir que as caixas sejam usadas \emph{em ordem}:
\[
y_i\ge y_{i+1}\qquad(i=1,\dots,n-1).
\]
Isso não muda o ótimo (sempre dá para renumerar as caixas usadas como $1,2,\dots,k$), mas elimina as cópias. No
experimento do notebook com 14 itens, sem a quebra de simetria o CBC bateu o limite de 30 s; com ela, resolveu em
0,15 s (ambos com 6 caixas). Os tempos variam com máquina e solver, mas a tendência é clara.

\begin{atencao}[erros típicos do Bloco 7]
\begin{itemize}
  \item Escrever a capacidade como $\sum_jw_jx_{ij}\le b$ (sem $y_i$): nada liga $y$ aos itens, e o solver zera
  todos os $y_i$, devolvendo ``0 caixas''.
  \item Usar $\sum_ix_{ij}\le1$ no lugar de $=1$: o modelo ``descarta'' itens e usa 0 caixas.
  \item Com \texttt{timeLimit}, o PuLP pode reportar ``Optimal'' para uma solução que é apenas viável. Desconfie de
  tempos que coincidem com o limite.
\end{itemize}
\end{atencao}

\chapter{Quadro-resumo: as receitas de linearização}
\label{cap:resumo}

\begin{center}
\small
\begin{tabularx}{\linewidth}{>{\raggedright\arraybackslash}p{4.6cm}l>{\raggedright\arraybackslash}X}
\toprule
Quero dizer\ldots & Variáveis & Restrição(ões) linear(es)\\
\midrule
$x>0\Rightarrow y=1$ (custo fixo, ativação) & $y\in\B$ & $x\le My$\\
$y=1\Rightarrow f(x)\le0$ & $y$ & $f(x)\le M(1-y)$\\
$y=1\Rightarrow x\ge m$ (lote mínimo) & $y$ & $x\ge m\,y$\\
$x\in\{0\}\cup[\ell,u]$ (semicontínua) & $y$ & $\ell y\le x\le uy$\\
$f(x)\le0$ \textbf{ou} $g(x)\le0$ & $y$ & $f\le M_1(1-y)$ \ e \ $g\le M_2\,y$\\
tarefas $j,k$ não se sobrepõem & $y_{jk}$ & $s_j+p_j\le s_k+M(1-y_{jk})$ \ e \ $s_k+p_k\le s_j+My_{jk}$\\
$P_1\Rightarrow P_2$ & $\delta_i$ & $\delta_1\le\delta_2$\\
$P_1\wedge P_2\Rightarrow P_3$ & $\delta_i$ & $\delta_1+\delta_2-\delta_3\le1$\\
$P_1\wedge\dots\wedge P_k\Rightarrow P_{k+1}\vee\dots\vee P_n$ & $\delta_i$ & $\sum_{i\le k}\delta_i-\sum_{i>k}\delta_i\le k-1$\\
no máximo / exatamente / pelo menos $k$ de $n$ & $\delta_i$ & $\sum\delta_i\le k$ \ / \ $=k$ \ / \ $\ge k$\\
$x\in\{v_1,\dots,v_k\}$ & $y_i$ & $x=\sum v_iy_i$, \ $\sum y_i=1$ \ (e $\phi(x)=\sum\phi(v_i)y_i$)\\
$\delta\leftrightarrow f(x)\le0$ & $\delta$ & $f\le M(1-\delta)$ \ e \ $\varepsilon-f-(\varepsilon-m)\delta\le0$\\
$P_1\xor P_2$ (com $\delta_i\leftrightarrow$ restrições) & $\delta_1,\delta_2$ & indicadores de cada uma \textbf{+} $\delta_1+\delta_2=1$\\
$\delta_3=\delta_1\delta_2$ & $\delta_3$ & $\delta_3\le\delta_1$, $\delta_3\le\delta_2$ (ou $-\delta_1-\delta_2+2\delta_3\le0$), $\delta_1+\delta_2-\delta_3\le1$\\
$\delta_3=\delta_1\vee\delta_2$ & $\delta_3$ & $\delta_3\ge\delta_1$, \ $\delta_3\ge\delta_2$, \ $\delta_3\le\delta_1+\delta_2$\\
$u=x\,\delta$ ($0\le x\le U$) & $u\ge0$ & $u\le U\delta$, \ $u\le x$, \ $u\ge x-U(1-\delta)$\\
\bottomrule
\end{tabularx}
\end{center}

\section*{Boas práticas (que voltam em \emph{branch-and-bound} e planos de corte)}
\begin{enumerate}
  \item \textbf{$M$ o menor possível.} Calcule-o com os limites das variáveis ou com uma pequena PL. $M$ grande
  enfraquece a relaxação e gera instabilidade numérica; $M$ pequeno demais deixa o modelo errado.
  \item \textbf{Formulações mais fortes} (desagregadas) podem valer mais que formulações ``enxutas''.
  \item \textbf{Simetria atrapalha.} Quebre-a com restrições de ordenação ($y_i\ge y_{i+1}$).
  \item \textbf{Modelo certo $\neq$ modelo bom.} Compare as \textbf{relaxações lineares} das alternativas.
  \item \textbf{Teste os casos.} Para cada restrição com binária, substitua $0$ e $1$ e leia o que sobra. Se
  alguma leitura não fizer sentido, a restrição está errada.
\end{enumerate}

\section*{Checklist antes de rodar o solver}
\begin{itemize}
  \item[$\square$] Cada variável tem o tipo certo (\py{Continuous}, \py{Integer}, \py{Binary}) e \py{lowBound=0} quando
  não pode ser negativa?
  \item[$\square$] O objetivo inclui \emph{todos} os custos/lucros, inclusive os fixos ($s_iy_i$)?
  \item[$\square$] Toda binária está ``amarrada'' a algo? (Uma binária que não aparece em nenhuma restrição o solver
  escolhe livremente.)
  \item[$\square$] Usei $=1$ onde o enunciado diz ``exatamente um'' e $\ge1$ onde diz ``pelo menos um''?
  \item[$\square$] Nenhum \py{<}, \py{>}, \py{!=} nem produto de variáveis?
  \item[$\square$] Os nomes das restrições são únicos (índice na f-string)?
  \item[$\square$] O resultado faz sentido? Confira a solução contra o enunciado, restrição por restrição.
\end{itemize}

\section*{Ponte para a lista de prática}
Os exercícios da lista (\emph{Pratica\_Exerc\_\_Programacao\_Inteira}) usam as técnicas deste notebook com problemas
diferentes:
\begin{center}
\begin{tabularx}{\linewidth}{lX}
\toprule
Exercício & Técnica(s)\\
\midrule
1 -- Artesão & Bloco 1: lote mínimo \emph{se} produzir; custo fixo\\
2 -- Máquinas paralelas & Blocos 6--7: alocação de tarefas a máquinas com objetivo \emph{min-max}\\
3 -- Tabuleiro $4\times4$ & lógica + variável inteira auxiliar (paridade: $\sum=2k$)\\
4 -- \emph{Flowshop} permutacional & Bloco 2: disjuntivas de ordenação (\emph{big-M})\\
\bottomrule
\end{tabularx}
\end{center}


\appendix
\chapter{Respostas dos exercícios de fixação}
\label{ap:respostas}

\section*{Capítulo 1}
\begin{description}
  \item[1.1] (a) PB (só binárias); (b) PLI (quantidades inteiras e moldes binários, sem contínuas); (c) PLIM
  (contínuas e uma binária).
  \item[1.2] (a) contínua basta; (b) inteira; (c) binária; (d) contínua basta.
  \item[1.3] A PL continua valendo 1000 ($o=5/3$). Inteiras: $o=0,v=5$ custa 1500; $o=1,v=2$ custa 1200; $o=2,v=0$
  custa 1200. O ótimo é 1200 (duas soluções), e arredondar $o=1{,}67\to2$ agora \emph{acerta}. Moral: às vezes
  arredondar funciona, mas não há garantia.
\end{description}

\section*{Capítulo 2}
\begin{description}
  \item[2.1] $50\cdot3+31\cdot3=243\le250$ e $3\cdot3-2\cdot3=3\ge-4$: viável, valor $3+1{,}92=4{,}92$. Não é
  ``vizinho'' porque os vizinhos do ótimo da PL têm $x_1\in\{1,2\}$ e $x_2\in\{4,5\}$.
  \item[2.2] A gulosa pega o item 1 (peso 3) e depois o item 2 (peso 2): peso 5, valor 19. A enumeração confirma que
  $(1,1,0)$, com valor 19, é ótimo. A gulosa acerta neste caso, por sorte.
  \item[2.3] Quarta-feira ($j=3$) é coberta por quem começou em qua, ter, seg, dom e sáb:
  $x_0+x_1+x_2+x_3+x_6\ge19$.
\end{description}

\section*{Capítulo 3}
\begin{description}
  \item[3.1] Em minimização, a relaxação é limite inferior e a incumbente é limite superior: o ótimo está em
  $[40,50]$. Então 38 não pode, 45 pode e 52 não pode (já temos uma solução de valor 50). Gap $=(50-40)/50=20\%$.
  \item[3.2] Essa solução é \textbf{ótima} para o problema binário.
  \item[3.3] A relaxação desse nó já é inteira ($x_1=1$), com valor $12\le13$: não há o que ramificar e ela é pior
  que a incumbente. O nó é podado.
\end{description}

\section*{Capítulo 4}
\begin{description}
  \item[4.1] $10y\le x\le50y$, \ $y\in\B$.
  \item[4.2] Queremos $\delta=0\Rightarrow f\ge0$ com $f=x_1-x_2-2$, cujo mínimo é $0-8-2=-10$. Restrição:
  $x_1-x_2-2\ge-10\delta$. Teste: $\delta=0$ dá $x_1-x_2\ge2$ \checkmark; \ $\delta=1$ dá $x_1-x_2\ge-8$, que sempre
  vale \checkmark.
  \item[4.3] $P_1\Rightarrow\neg P_2\equiv\neg P_1\vee\neg P_2$, que vira $(1-\delta_1)+(1-\delta_2)\ge1$, ou seja,
  $\delta_1+\delta_2\le1$. Só a linha $(1,1)$ é proibida \checkmark.
  \item[4.4] $f_i(x)\le M_i(1-\delta_i)$ para $i=1,2,3$, e $\delta_1+\delta_2+\delta_3\ge2$.
  \item[4.5] (1) A relaxação fica fraca, e o \emph{branch-and-bound} demora muito mais; (2) problemas numéricos: uma
  binária ``quase zero'' (dentro da tolerância de integralidade) multiplicada por $10^9$ libera valores grandes de
  $x$, o fenômeno que a Gurobi chama de \emph{trickle flow} (Seção~\ref{sec:escolherM}).
\end{description}

\section*{Capítulo 5}
\begin{description}
  \item[5.1] Veja a função \py{modelo_wolsey} no Capítulo~\ref{cap:bloco0}.
\end{description}

\chapter{Gabarito de valores ótimos do notebook}

Valores do dicionário \py{ESPERADO} do notebook, para conferir sem rodar a verificação.
\begin{center}
\begin{tabular}{lll}
\toprule
Chave & Valor & Onde\\
\midrule
\texttt{wolsey\_lp} / \texttt{wolsey\_ip} & 5,0984 / 5 & Bloco 0\\
\texttt{fabrica} / \texttt{fabrica\_regra} & 1875 / 1375 & Bloco 1\\
\texttt{M1} / \texttt{M2} / \texttt{disjuncao} & 120 / 20 / 100 & Bloco 2A\\
\texttt{1maq} / \texttt{1maq\_release} & 127 / 135 & Bloco 2B\\
\texttt{investimentos} & 21 & Bloco 3B\\
\texttt{inv\_rel\_agregada} / \texttt{inv\_rel\_desagregada} & 23,714 / 22,5 & Bloco 3B\\
\texttt{tubos} & 41\,100 & Bloco 4\\
\texttt{ressonancia} & 39 & Bloco 5A\\
\texttt{desconto\_50} / \texttt{\_75} / \texttt{\_90} / \texttt{\_130} & 500 / 750 / 800 / 1040 & Bloco 5B\\
\texttt{mochila} / \texttt{mochila\_extras} & 142 / 135 & Bloco 6\\
\texttt{binpacking} & 6 & Bloco 7\\
\bottomrule
\end{tabular}
\end{center}


\begin{thebibliography}{99}
\addcontentsline{toc}{chapter}{Referências}

\bibitem{notebook}
Moreira, A.~C.~B.; Moreira, M.~C.~O.
\emph{Aula\_Programacao\_Inteira\_Modelagem.ipynb} -- GCC118 Programação Matemática. DMM/DCC, Universidade Federal de
Lavras, 2026. (Notebook de aula; base desta apostila.)

\bibitem{slides}
Toledo, F.~M.~B.; Costa, A.~M.
\emph{Programação inteira -- Modelagem} (slides de aula, ICMC/USP), com slides do Prof.\ Celso Carneiro Ribeiro (UFF).
Material de apoio de GCC118.

\bibitem{wolsey1998}
Wolsey, L.~A. \emph{Integer Programming}. New York: Wiley, 1998.

\bibitem{arenales2007}
Arenales, M.; Armentano, V.; Morabito, R.; Yanasse, H. \emph{Pesquisa Operacional}. Rio de Janeiro: Elsevier/Campus,
2007.

\bibitem{sierksma1996}
Sierksma, G. \emph{Linear and Integer Programming: Theory and Practice}. New York: Marcel Dekker, 1996.

\bibitem{amp}
Bradley, S.~P.; Hax, A.~C.; Magnanti, T.~L. \emph{Applied Mathematical Programming}. Reading: Addison-Wesley, 1977.
Cap.~9: \emph{Integer Programming}. Disponível em: \url{http://web.mit.edu/15.053/www/AMP-Chapter-09.pdf}. Acesso em:
out.\ 2026.

\bibitem{landdoig1960}
Land, A.~H.; Doig, A.~G. An automatic method of solving discrete programming problems. \emph{Econometrica},
v.~28, n.~3, p.~497--520, 1960.

\bibitem{informsland}
INFORMS. \emph{Ailsa H. Land} (perfil biográfico, History of O.R. Excellence). Disponível em:
\url{https://www.informs.org/Explore/History-of-O.R.-Excellence/Biographical-Profiles/Land-Ailsa-H}. Acesso em:
out.\ 2026.

\bibitem{gomory1958}
Gomory, R.~E. Outline of an algorithm for integer solutions to linear programs. \emph{Bulletin of the American
Mathematical Society}, v.~64, n.~5, p.~275--278, 1958.

\bibitem{dakin1965}
Dakin, R.~J. A tree-search algorithm for mixed integer programming problems. \emph{The Computer Journal}, v.~8,
n.~3, p.~250--255, 1965.

\bibitem{bixby2012}
Bixby, R.~E. A brief history of linear and mixed-integer programming computation. \emph{Documenta Mathematica},
Extra Volume ISMP (Optimization Stories), p.~107--121, 2012.

\bibitem{koch2022}
Koch, T.; Berthold, T.; Pedersen, J.; Vanaret, C. Progress in mathematical programming solvers from 2001 to 2020.
\emph{EURO Journal on Computational Optimization}, v.~10, 100031, 2022. Preprint: \url{https://arxiv.org/abs/2206.09787}.

\bibitem{vielma2015}
Vielma, J.~P. Mixed integer linear programming formulation techniques. \emph{SIAM Review}, v.~57, n.~1, p.~3--57,
2015. Versão aberta: \url{https://dspace.mit.edu/handle/1721.1/96480}.

\bibitem{browndell2007}
Brown, G.~G.; Dell, R.~F. Formulating integer linear programs: a rogues' gallery. \emph{INFORMS Transactions on
Education}, v.~7, n.~2, p.~153--159, 2007.

\bibitem{williams2013}
Williams, H.~P. \emph{Model Building in Mathematical Programming}. 5.~ed. Chichester: Wiley, 2013.

\bibitem{williams2009}
Williams, H.~P. \emph{Logic and Integer Programming}. New York: Springer, 2009.

\bibitem{ravelo}
Ravelo, S.~V. \emph{Programação linear inteira} (aula 22, MC558 -- Projeto e Análise de Algoritmos II). Instituto de
Computação, Unicamp. Disponível em: \url{https://ic.unicamp.br/~santiago/assets/mc558/MC558___22.pdf}. Acesso em:
out.\ 2026.

\bibitem{gurobitrickle}
Gurobi Optimization. \emph{Why does Gurobi report my model is infeasible when it has a feasible solution?} Gurobi
Help Center. Disponível em: \url{https://support.gurobi.com/hc/en-us/articles/16565310050833}. Acesso em: out.\ 2026.

\bibitem{gurobibigm}
Gurobi Optimization. \emph{Dealing with big-M constraints} (Gurobi Reference Manual, guia de questões numéricas).
Disponível em: \url{https://gurobi.com/documentation/9.0/refman/num_dealing_with_big_m_con.html}. Acesso em: out.\ 2026.

\bibitem{amplnum}
AMPL Optimization. \emph{Numerical accuracy} (documentação do AMPL MP). Disponível em:
\url{https://mp.ampl.com/modeling-numerics.html}. Acesso em: out.\ 2026.

\bibitem{smith1956}
Smith, W.~E. Various optimizers for single-stage production. \emph{Naval Research Logistics Quarterly}, v.~3,
n.~1--2, p.~59--66, 1956.

\bibitem{dantzig1957}
Dantzig, G.~B. Discrete-variable extremum problems. \emph{Operations Research}, v.~5, n.~2, p.~266--288, 1957.

\bibitem{martellototh1990}
Martello, S.; Toth, P. \emph{Knapsack Problems: Algorithms and Computer Implementations}. Chichester: Wiley, 1990.

\bibitem{kantorovich1960}
Kantorovich, L.~V. Mathematical methods of organizing and planning production. \emph{Management Science}, v.~6,
n.~4, p.~366--422, 1960. (Tradução do original russo de 1939.)

\bibitem{gilmoregomory1961}
Gilmore, P.~C.; Gomory, R.~E. A linear programming approach to the cutting-stock problem. \emph{Operations
Research}, v.~9, n.~6, p.~849--859, 1961.

\bibitem{carvalho2002}
Valério de Carvalho, J.~M. LP models for bin packing and cutting stock problems. \emph{European Journal of
Operational Research}, v.~141, n.~2, p.~253--273, 2002.

\bibitem{delorme2016}
Delorme, M.; Iori, M.; Martello, S. Bin packing and cutting stock problems: mathematical models and exact
algorithms. \emph{European Journal of Operational Research}, v.~255, n.~1, p.~1--20, 2016.

\bibitem{ocw15083}
Bertsimas, D.; Schulz, A. \emph{15.083J Integer Programming and Combinatorial Optimization} (Fall 2009). MIT
OpenCourseWare. Disponível em:
\url{https://ocw.mit.edu/courses/15-083j-integer-programming-and-combinatorial-optimization-fall-2009/}. Acesso em:
out.\ 2026.

\bibitem{pulp2011}
Mitchell, S.; O'Sullivan, M.; Dunning, I. \emph{PuLP: a linear programming toolkit for Python}. The University of
Auckland, 2011. Código e documentação: \url{https://github.com/coin-or/PuLP}.

\end{thebibliography}


\end{document}
